% ============================================
%  校园二手拍卖与即时通讯平台 (CAIP)
%  需求与功能设计规格书
%  文件编号: CAIP-FD-001
%  生成日期: 2026-06-11
%  使用 XeLaTeX 编译
% ============================================

\documentclass[12pt,a4paper]{ctexart}

% ==================== 页面设置 ====================
\usepackage[top=2.54cm, bottom=2.54cm, left=3.17cm, right=3.17cm]{geometry}
\usepackage{setspace}
\onehalfspacing

% ==================== 字体设置 ====================
% 本地编译 (Windows): 取消下面5行注释, 注释掉 Overleaf 字体部分
% \setCJKmainfont{SimSun}[AutoFakeBold=2.5]  % 宋体 正文
% \setCJKsansfont{SimHei}[AutoFakeBold=2.5]  % 黑体 标题
% \setCJKmonofont{SimSun}
% \setmainfont{Times New Roman}

% Overleaf 编译: 使用 Fandol 字体 (TeX Live 自带)
\setCJKmainfont{FandolSong}[
  BoldFont=FandolHei,
  ItalicFont=FandolKai,
  AutoFakeBold=2.5
]
\setCJKsansfont{FandolHei}[AutoFakeBold=2.5]
\setCJKmonofont{FandolFang}
\setmainfont{TeX Gyre Termes}  % Times New Roman 的替代

% ==================== 常用宏包 ====================
\usepackage{graphicx}
\usepackage{float}
\usepackage{longtable}
\usepackage{booktabs}
\usepackage{array}
\usepackage{multirow}
\usepackage{hyperref}
\usepackage{fancyhdr}
\usepackage{titletoc}
\usepackage{caption}
\usepackage{enumitem}
\usepackage{xcolor}
\usepackage{colortbl}
\usepackage{fancyvrb}
\usepackage{amsmath}
\usepackage{amssymb}
\usepackage[ export ]{adjustbox}

% ==================== 代码块 ====================
\usepackage{listings}
\lstset{
    basicstyle=\ttfamily\small,
    breaklines=true,
    frame=single,
    numbers=none,
    aboveskip=1em,
    belowskip=1em,
    showstringspaces=false,
    tabsize=2
}

% ==================== 超链接设置 ====================
\hypersetup{
    colorlinks=true,
    linkcolor=black,
    filecolor=black,
    urlcolor=black,
    citecolor=black,
    pdfborder={0 0 0}
}

% ==================== 页眉页脚 ====================
\fancypagestyle{plain}{
    \fancyhf{}
    \fancyfoot[C]{\thepage}
    \renewcommand{\headrulewidth}{0pt}
}
\pagestyle{plain}

% ==================== 标题格式 ====================
\ctexset{
    section = {
        format = {\zihao{3}\bfseries\heiti\centering},
        beforeskip = 24pt,
        afterskip = 12pt
    },
    subsection = {
        format = {\zihao{4}\bfseries\heiti},
        beforeskip = 18pt,
        afterskip = 9pt
    },
    subsubsection = {
        format = {\zihao{-4}\bfseries\heiti},
        beforeskip = 12pt,
        afterskip = 6pt
    }
}

% ==================== 表格颜色 ====================
\definecolor{tableheader}{RGB}{217,225,242}

% ==================== 图表编号 ====================
\numberwithin{figure}{section}
\numberwithin{table}{section}
\renewcommand{\figurename}{图}
\renewcommand{\tablename}{表}

% ==================== 目录格式 ====================
\renewcommand{\contentsname}{\centering\zihao{3}\heiti 目\quad 录}

\begin{document}

% #################### 封面 ####################
\begin{titlepage}
\thispagestyle{empty}
\begin{center}

\vspace*{2cm}

{\zihao{-3} 文件编号：CAIP-FD-001}

\vspace{2cm}

{\zihao{2}\heiti 项目名：}

\vspace{0.5cm}

{\zihao{2}\heiti 校园二手拍卖与即时通讯平台}

\vspace{0.3cm}

{\zihao{2}\heiti (Campus Auction \& IM Platform, CAIP)}

\vspace{2cm}

{\zihao{1}\bfseries\heiti 需求与功能设计规格书}

\vspace{3cm}

{
\zihao{4}
\begin{tabular}{rl}
\textbf{作者：} & 顾思豪、刘子彦、段睿智、罗圣杰、罗涵、欧阳熠 \\
& \\
\textbf{部门：} & 第9组 \\
& \\
\textbf{时间：} & 2026年6月 \\
\end{tabular}
}

\vspace{2cm}

{\zihao{3}\heiti 软件工程实验}

\vfill

{\zihao{5} 2026年春季学期 · 第13-16周}

\end{center}
\end{titlepage}

% #################### 改版履历 ####################
\newpage
\section*{\centering 改版履历}
\addcontentsline{toc}{section}{改版履历}

\vspace{1em}

\begin{center}
\renewcommand{\arraystretch}{1.5}
\begin{tabular}{|c|c|p{6cm}|c|c|c|p{3cm}|}
\hline
\rowcolor{tableheader}
\textbf{版数} & \textbf{年月日} & \textbf{改版内容} & \textbf{承认} & \textbf{查阅} & \textbf{作成} & \textbf{备注} \\
\hline
v1.0 & 2026-06-09 & 初稿作成 & --- & --- & 9组全员 & --- \\
\hline
\end{tabular}
\end{center}

\vspace{2em}

% ==================== 模板说明 (保留但以注释形式) ====================
\begin{quote}
{\kaishu\zihao{5}【说明】以下虚线部分为模板说明内容。虚线部分删除后填入记述内容。

以下对设计书中记述的项目进行说明。

目录中的项目原则上不能删除。如果没有记载事项的场合，要明确记述其原因。但是，由于开发部件的种类或者特性，必须删除目录中的项目时，要在过程定义中明确删除的位置和理由。

正文统一采用SimSun（宋体）五号字体。下同，不再重复说明。}
\end{quote}

% #################### 目录 ####################
\newpage
\tableofcontents
\newpage

% #################### 第1章 前言 ####################
\section{前言}

要理解本设计书，必须明确相关的知识。说明本文档的背景。

\subsection{输入文档}

明确编写本需求设计书时输入的文档。通常，需求设计书的输入文档如下所示：

\begin{itemize}
    \item 需求定义书
    \item 开发计划书
\end{itemize}

\begin{center}
\renewcommand{\arraystretch}{1.3}
\begin{tabular}{|c|c|c|c|}
\hline
\rowcolor{tableheader}
\textbf{编号} & \textbf{文档名称} & \textbf{版本} & \textbf{来源} \\
\hline
CAIP-PLAN-001 & 校园二手拍卖与即时通讯平台 开发计划书 & v1.1 & 9组 \\
\hline
课程-001 & 软件工程实验总要求 & 2026 & 课程 \\
\hline
TEMPLATE-001 & 02\_FD 功能和概要设计说明书 文档模板 & --- & 课程模板 \\
\hline
\end{tabular}
\end{center}

\subsection{相关文档}

明确与本需求设计书中规定功能相关联的其他文档。通常，本需求设计书的输入式样书如下所示：

\begin{itemize}
    \item 与本功能相关功能的概要需求定义书
    \item 开发计划书
\end{itemize}

\begin{center}
\renewcommand{\arraystretch}{1.3}
\begin{tabular}{|c|c|c|}
\hline
\rowcolor{tableheader}
\textbf{编号} & \textbf{文档名称} & \textbf{说明} \\
\hline
CAIP-PLAN-001 & 开发计划书 & 项目规划、需求规格、技术选型 \\
\hline
CAIP-TEST-001 & 测试设计及结果报告书 & 测试策略与测试结果（后续交付） \\
\hline
CAIP-SUM-001 & 项目总结报告书 & 项目回顾与经验总结（后续交付） \\
\hline
\end{tabular}
\end{center}

\subsection{用语/术语定义}

明确以下记述中所使用的用语/术语定义。在本章中主要明确以下用语/术语：

\begin{itemize}
    \item 向用户公开的用语/术语
    \item 本次开发特有的用语/术语
    \item 用语/术语本身并不特殊，但在本次开发中却有特殊含义的用语/术语
\end{itemize}

\textbf{单独做成用语集时，请把链接粘贴到本文档中。}

\begin{center}
\renewcommand{\arraystretch}{1.3}
\begin{tabular}{|p{3cm}|p{11cm}|}
\hline
\rowcolor{tableheader}
\textbf{术语} & \textbf{定义} \\
\hline
CAIP & 本系统英文简称，校园二手拍卖与即时通讯平台 \\
\hline
英式拍卖 & 出价从低到高的升价拍卖，在截止时间内价高者得 \\
\hline
一口价 & 卖家设定固定价格，买家直接购买，无需竞价 \\
\hline
流拍 & 拍卖结束时无人出价，商品状态变为EXPIRED \\
\hline
STOMP & Simple Text Oriented Messaging Protocol，面向文本的消息协议，运行于WebSocket之上 \\
\hline
JWT & 无状态认证令牌，包含用户ID、角色、过期时间，使用HMAC-SHA256签名 \\
\hline
MyBatis-Plus & Java持久层框架，提供Lambda表达式查询、自动分页 \\
\hline
Flyway & 数据库版本迁移工具，通过SQL脚本管理数据库Schema变更 \\
\hline
IM & 即时通讯，本系统中买卖双方一对一私聊功能 \\
\hline
DFD & 数据流图，描述系统中数据的流动方向和处理逻辑 \\
\hline
RBAC & 基于角色的访问控制，通过角色控制用户权限 \\
\hline
HikariCP & 高性能JDBC数据库连接池 \\
\hline
ORM & 对象关系映射，使用MyBatis-Plus实现 \\
\hline
\end{tabular}
\end{center}

% #################### 第2章 开发概要 ####################
\section{开发概要}

\subsection{开发的目的和背景}

描述本设计书中所记载功能的开发目的。注意考虑经济可行性的同时，要考虑到是否符合相应的法律法规以及道德的准则。

\textbf{目的：}

构建面向高校学生的校园二手拍卖与即时通讯一体化平台。以实时拍卖模式实现公允的价格发现，以内置即时通讯(IM)替代第三方通信工具保护用户隐私，以校园身份认证（学号/工号）构建可信交易环境，以信誉评价体系降低交易风险。

\textbf{背景：}

高校内每年产生大量二手物品交易需求。现有交易渠道存在严重问题：QQ群/微信群消息刷屏，商品信息难以检索和追溯；买卖双方须添加微信或QQ沟通，隐私泄露风险高；缺乏拍卖等趣味性交易模式；先款后货缺乏信任保障。

本系统作为软件工程实验课程项目，完整实践需求分析→系统设计→编码实现→测试交付的全流程，符合课程对6人团队协作完成设计、编码、测试的要求。

\textbf{经济可行性：}

全部技术栈（Spring Boot、Vue 3、MySQL、Redis）均为开源免费软件，无需支付授权费用。部署于学校服务器或低成本云主机，运营成本低。

\textbf{法律法规与道德准则：}

\begin{itemize}
    \item 用户数据（学号、邮箱、手机号）仅用于平台身份认证，不向第三方提供
    \item IM聊天内容敏感词过滤，防止传播违法违规信息
    \item 二手交易限于合法物品，禁止交易违禁品（商品描述审核机制）
    \item 用户评价须基于真实交易，禁止刷评价
\end{itemize}

\subsection{工作\&运行环境}

明确本产品相关的工作环境（HW/SW）。

\begin{itemize}
    \item 作为HW环境，明确使要开发功能能够运行的必要条件（CPU、内存、机器种类等）
    \item 作为SW环境，明确使要开发功能能够运行的OS、相关SW等
    \item 系统部署图
\end{itemize}

\subsubsection{服务器端运行环境}

\textbf{■ HW（硬件环境）}

\begin{itemize}
    \item CPU：最低2核（x86\_64架构），推荐4核以上
    \item 内存：最低4GB，推荐8GB以上
    \item 磁盘：最低20GB（不含用户上传图片），推荐40GB SSD
    \item 网络：最低100Mbps，推荐1Gbps
    \item 机器种类：云主机 / 虚拟机 / 物理服务器均可，无特殊硬件要求
\end{itemize}

阿里云ECS（2核4G，40GB SSD）/ 腾讯云CVM / 华为云ECS / 校内虚拟化平台虚拟机

\textbf{■ SW（软件环境）}

\begin{itemize}
    \item OS：Ubuntu 20.04 LTS 64bit 或 CentOS 7+ 64bit（Linux发行版）
    \item Java运行时：OpenJDK 21（Spring Boot 3.4 要求 JDK 17+）
    \item 数据库：MySQL 8.0.36
    \item 缓存：Redis 7.2
    \item Web服务器：Nginx 1.24+（反向代理 + 静态资源托管）
    \item 容器运行时：Docker Engine 24+ + Docker Compose v2（可选，用于容器化部署）
\end{itemize}

使用Docker Compose统一管理所有服务，避免环境差异导致的问题。

\subsubsection{客户端运行环境}

\textbf{■ HW（硬件环境）}

\begin{itemize}
    \item CPU：最低1核，推荐2核以上
    \item 内存：最低2GB，推荐4GB以上
    \item 屏幕分辨率：最低375px宽（手机竖屏），推荐1920×1080（桌面）
    \item 网络：最低1Mbps，推荐5Mbps以上
    \item 机器种类：PC、笔记本电脑、智能手机、平板电脑均可
\end{itemize}

\textbf{■ SW（软件环境）}

\begin{itemize}
    \item OS：Windows 10+ / macOS 12+ / Android 10+ / iOS 15+ / 任意Linux发行版
    \item 对应Web浏览器（明确最低版本）：
    \begin{itemize}
        \item 桌面端：Google Chrome 120+、Microsoft Edge 120+、Mozilla Firefox 120+
        \item 移动端：iOS Safari 16+、Android Chrome 120+
    \end{itemize}
\end{itemize}

使用Chrome或Edge浏览器以获得最佳WebSocket连接稳定性。

\begin{quote}
\textbf{注意}：本系统为Web应用程序，客户端仅需浏览器，无需安装任何客户端软件或浏览器插件。
\end{quote}

\subsubsection{系统部署图}

\begin{figure}[H]
\centering
{\small
\begin{verbatim}
┌─────────────────────────────┐    ┌───────────────────────────────────────────┐
│          用户端              │    │         Docker Compose 环境               │
│                             │    │         (Ubuntu 20.04)                     │
│  ┌───────────────────────┐  │    │                                            │
│  │ 浏览器                 │  │    │  ┌──────────┐  ┌────────┐  ┌────────┐   │
│  │ Chrome/Edge/Safari    │  │    │  │  Nginx   │  │Spring  │  │ MySQL  │   │
│  │ 需支持WebSocket        │  │    │  │  :443    │  │Boot    │  │ 8.0    │   │
│  └───────────────────────┘  │    │  │  HTTPS/  │  │:8080   │  │ :3306  │   │
│                             │    │  │  WSS     │  │REST API│  │ 数据持久│   │
└─────────────────────────────┘    │  │  反向代理 │  │+WS     │  │ 化卷   │   │
         │                         │  └────┬─────┘  └───┬────┘  └───┬────┘   │
         │ HTTPS :443              │       │              │            │       │
         │ WSS :443                │       │    反向代理   │   JDBC     │       │
         │                         │       └──────────────┤            │       │
         └─────────────────────────┘                      │            │       │
                                                          │  ┌─────────┘       │
                                                          │  │                 │
                                                          │  │  ┌──────────┐   │
                                                          │  └──┤  Redis 7 │   │
                                                          │     │  :6379   │   │
                                                          │     │  RDB持久化 │   │
                                                          │     └──────────┘   │
                                                          │  Lettuce           │
                                                          └────────────────────┘
\end{verbatim}
}
\caption{系统部署图}
\label{fig:deploy}
\end{figure}

\subsection{系统构成}

在系统构成中明确以下内容：

\begin{itemize}
    \item 明确本产品预期的系统构成
    \item 明确所开发功能的系统整体位置
    \item 明确与其他运行程序的关联关系
\end{itemize}

\textbf{可使用UML配置图表示系统构成。}

本系统采用前后端分离 + 三层架构的体系结构。前端Vue 3单页应用由Nginx托管静态文件。后端Spring Boot应用按6个功能模块垂直拆分，公共横切关注点（Security/JWT/Exception/Logging）抽取为common模块。数据访问层使用MyBatis-Plus操作MySQL，Spring Data Redis操作Redis缓存。实时通信层使用Spring WebSocket + STOMP协议内嵌于业务模块中。

\textbf{【功能模块划分】}

\begin{itemize}
    \item \textbf{M1 --- 用户与认证模块}：注册登录、JWT鉴权、RBAC权限、个人信息管理
    \item \textbf{M2 --- 商品发布与搜索模块}：发布（一口价/拍卖）、多图上传、全文搜索、分类筛选、收藏
    \item \textbf{M3 --- 实时拍卖引擎模块}：英式升价拍卖、Redis Lua原子出价、STOMP实时推送、定时判定成交
    \item \textbf{M4 --- 即时通讯模块(IM)}：一对一私聊、文字+图片、WebSocket实时收发、已读标记
    \item \textbf{M5 --- 订单与交易模块}：订单状态机、模拟支付、发货确认收货
    \item \textbf{M6 --- 管理后台与统计模块}：用户审核、商品管理、数据统计ECharts看板
\end{itemize}

\begin{figure}[H]
\centering
{\small
\begin{verbatim}
┌─────────────────────────────────────────────────────────────┐
│                 表示层 Presentation                          │
│            Vue 3 + Element Plus / Nginx静态托管               │
└──────────────────────────┬──────────────────────────────────┘
                           │
┌──────────────────────────┴──────────────────────────────────┐
│            业务逻辑层 Business Logic (Spring Boot 3.4)        │
│  ┌─────────────────┐  ┌──────────────┐  ┌──────────────┐   │
│  │ M1-M6 六大模块   │  │ 公共模块      │  │ WebSocket    │   │
│  │                 │  │ security/    │  │ STOMP        │   │
│  │                 │  │ config/      │  │              │   │
│  │                 │  │ exception    │  │              │   │
│  └────────┬────────┘  └──────────────┘  └──────────────┘   │
└───────────┼────────────────────────────────────────────────┘
            │
┌───────────┴────────────────────────────────────────────────┐
│                 数据访问层 Data Access                       │
│     MyBatis-Plus      Spring Data Redis      Flyway迁移     │
└───────────┬────────────────┬──────────────────────────────┘
            │                │
┌───────────┴────────┐  ┌───┴──────────┐  ┌─────────────────┐
│    MySQL 8.0       │  │   Redis 7    │  │  本地文件系统    │
│    (数据存储)       │  │   (缓存)     │  │  (图片存储)     │
└────────────────────┘  └──────────────┘  └─────────────────┘
\end{verbatim}
}
\caption{系统体系结构图}
\label{fig:architecture}
\end{figure}

\subsection{功能整体图}

可以使用模块层次方框图。也可以使用用例图。

除了图以外，还要记述本次开发功能整体的UML中不能记述的事项或者补充事项。

\subsubsection{模块层次方框图}

\begin{figure}[H]
\centering
{\small
\begin{verbatim}
校园二手拍卖与即时通讯平台 (CAIP)
│
├── M1 用户与认证模块 ─── 注册 / 登录 / JWT / RBAC
│   └── 子功能：注册 登录 Token刷新 个人信息 用户主页 管理员审核封禁
│
├── M2 商品发布与搜索模块 ─── 发布 / 搜索 / 筛选 / 收藏 / 图片
│   └── 子功能：发布商品 商品列表 详情 全文搜索 多条件筛选 收藏 编辑
│
├── M3 实时拍卖引擎模块 ─── 出价 / Redis Lua / STOMP推送 / 成交
│   └── 子功能：出价 实时推送 倒计时 定时成交判定 出价历史 当前价查询
│
├── M4 即时通讯模块 ─── 私聊 / WebSocket / 已读 / 文字图片
│   └── 子功能：会话列表 创建会话 消息历史 发送消息 已读标记
│
├── M5 订单与交易模块 ─── 状态机 / 模拟支付 / 发货收货
│   └── 子功能：创建订单 模拟支付 标记发货 确认收货 取消订单
│
├── M6 管理后台与统计模块 ─── 审核 / 封禁 / ECharts看板
│   └── 子功能：用户审核 封禁解封 商品审核 分类管理 数据概览 趋势图
│
└── 公共横切模块 (Common)
     Spring Security配置 / JWT过滤器 / 全局异常处理器 / AOP操作日志 / 统一响应格式
\end{verbatim}
}
\caption{模块层次方框图}
\label{fig:module-hierarchy}
\end{figure}

\subsubsection{系统用例图}

\begin{figure}[H]
\centering
{\small
\begin{verbatim}
                        CAIP 系统边界
  ┌──────────────────────────────────────────────────────┐
  │                                                      │
  │  ┌─────────────┐  ┌─────────────┐  ┌──────────────┐ │
  │  │  注册账号    │  │  浏览/搜索   │  │  登录系统    │ │
  │  └─────────────┘  │  商品       │  └──────────────┘ │
  │                   └─────────────┘                    │
  │  ┌─────────────┐  ┌─────────────┐  ┌──────────────┐ │
  │  │  发布商品    │  │  参与竞拍    │  │  IM私聊     │ │
  │  └─────────────┘  └─────────────┘  └──────────────┘ │
  │                                                      │
  │  ┌─────────────┐  ┌─────────────┐  ┌──────────────┐ │
  │  │ 下单/支付/   │  │  评价交易    │  │  管理个人    │ │
  │  │ 收货        │  └─────────────┘  │  信息        │ │
  │  └─────────────┘                   └──────────────┘ │
  │                                                      │
  │  ┌─────────────┐  ┌─────────────┐                   │
  │  │ 审核用户/    │  │  查看数据    │                   │
  │  │ 商品        │  │  统计       │                   │
  │  └─────────────┘  └─────────────┘                   │
  └──────────────────────────────────────────────────────┘
     ▲            ▲                ▲
     │            │                │
  ┌──┴──┐    ┌───┴───┐      ┌───┴───┐
  │访客  │    │注册用户│      │管理员  │
  └─────┘    └───────┘      └───────┘
\end{verbatim}
}
\caption{系统用例图}
\label{fig:use-case}
\end{figure}

\subsubsection{UML图中无法记述的补充事项}

模块层次方框图和用例图表达了系统的静态功能分解和参与者-用例关系，但以下动态行为和设计约束无法在图中充分体现，在此补充说明。

\textbf{(1) 模块间交互关系（非方框图的层次关系）}

方框图仅表达了模块的层次分解，未表达模块间的动态调用关系。实际上6个模块之间存在紧密的业务协作：

\begin{itemize}
    \item \textbf{M3拍卖引擎 → M5订单交易}：拍卖成交后，M3自动调用M5的 \verb|OrderService.createOrder()| 接口生成订单。这是最关键的跨模块调用链，涉及Redis读取winner → products表更新 → orders表插入的事务一致性
    \item \textbf{M3拍卖引擎 → M4即时通讯}：拍卖成交或流拍后，M3通过M4的WebSocket通道向买卖双方推送系统通知消息
    \item \textbf{M2商品发布 → M3拍卖引擎}：type='AUCTION'的商品发布后，M3读取auction\_start/auction\_end字段接管拍卖生命周期
    \item \textbf{M2商品发布 → M6管理后台}：新用户首次发布商品须经M6管理员审核通过后才公开
    \item \textbf{M1用户认证 → 全部模块}：所有模块的接口调用均须经过M1的JwtAuthenticationFilter进行Token校验
    \item \textbf{M4即时通讯 → M2商品浏览}：IM会话可通过product\_id关联到具体商品
\end{itemize}

\textbf{(2) 功能优先级（MVP增量交付策略）}

采用增量开发模型，功能按优先级分批交付：

\begin{itemize}
    \item \textbf{Sprint 1（核心功能，必须交付）}：M1用户注册登录(JWT) + M2商品发布浏览 + M3拍卖出价核心逻辑
    \item \textbf{Sprint 2（完整功能，必须交付）}：M4即时通讯 + M5订单状态机 + M6管理后台基本功能
    \item \textbf{扩展功能（选做，视进度）}：代理出价、校内摊位LBS地图、热门商品推荐
\end{itemize}

\textbf{(3) 功能间的数据依赖关系}

\begin{itemize}
    \item 拍卖出价(M3)依赖M2提供的商品起拍价(start\_price)和拍卖结束时间(auction\_end)
    \item 订单生成(M5)依赖M2提供的商品卖家信息(seller\_id)和M3提供的成交价(current\_price)与得主(winner\_id)
    \item IM会话(M4)可选关联M2的product\_id，形成"针对某商品的对话"的业务上下文
    \item 信誉评价(M6)依赖M5的订单完成状态(COMPLETED)作为评价资格的判断条件
\end{itemize}

\textbf{(4) 非功能约束在各模块中的体现}

\begin{itemize}
    \item \textbf{安全性}：所有模块接口（除登录注册外）均须JWT认证。拍卖出价须校验出价者≠卖家本人。订单操作须校验操作者为买家或卖家本人。管理后台须校验ADMIN角色
    \item \textbf{性能}：M3拍卖出价使用Redis Lua脚本保证原子性（<1ms延迟）。M2商品搜索使用MySQL ngram全文索引。M4 IM消息使用WebSocket直推
    \item \textbf{并发控制}：M3出价使用Redis Lua脚本原子操作。M2库存使用乐观锁。M5订单状态变更使用版本号乐观锁
    \item \textbf{数据一致性}：M3拍卖成交后更新product表+创建order表在同一数据库事务中完成
\end{itemize}

\textbf{(5) 用户旅程与模块切换路径}

\begin{itemize}
    \item \textbf{买家拍卖旅程}：浏览商品(M2)→查看详情(M2)→出价(M3)→实时收到推送(M3)→拍卖成交→查看订单(M5)→IM沟通细节(M4)→模拟支付(M5)→确认收货(M5)→评价(M6)
    \item \textbf{卖家发布旅程}：发布拍卖商品(M2)→等待审核(M6)→拍卖开始→收到出价推送(M3)→拍卖成交→查看订单(M5)→IM沟通(M4)→发货(M5)→确认买家收货→收到评价(M6)
    \item \textbf{管理员工作流}：登录→仪表盘概览(M6)→待审核用户列表(M6)→审核通过/驳回(M6)→待审核商品列表(M6)→通过/下架(M6)→查看交易趋势统计(M6)
\end{itemize}

\subsection{用户分类和特征}

使用表的方式，明确以下内容：

\begin{itemize}
    \item 本系统的用户对象
    \item 用户对象之间的关系（给出关系图）
    \item 用户和功能模块的对应关系
\end{itemize}

\subsubsection{用户一览表}

\begin{center}
\renewcommand{\arraystretch}{1.3}
\begin{tabular}{|c|c|c|p{6cm}|}
\hline
\rowcolor{tableheader}
\textbf{NO} & \textbf{用户} & \textbf{别名/角色标识} & \textbf{简要说明} \\
\hline
1 & 未登录访客 & ANONYMOUS & 未登录的浏览者，可查看公开商品信息和用户主页，但不能交易 \\
\hline
2 & 普通用户（买家/卖家） & USER & 在校学生或教职工，经学号/工号认证注册。同一账号可兼具买卖家身份 \\
\hline
3 & 管理员 & ADMIN & 平台运营维护者，由系统在数据库中直接指定（不开放注册），负责审核和监管 \\
\hline
\end{tabular}
\end{center}

\subsubsection{各用户角色的职责}

\textbf{1. 未登录访客（ANONYMOUS）}

通过浏览器访问平台首页，可以浏览热门商品和推荐内容、搜索商品、查看商品详情。可以查看任意用户的公开主页。但无法发布商品、出价、私聊或下单。

\begin{itemize}
    \item 浏览商品Feed流（首页推荐 + 分类浏览）
    \item 搜索商品（关键词 + 分类 + 价格区间筛选）
    \item 查看商品详情（图片轮播、商品信息、出价记录）
    \item 查看用户公开主页（信誉分、评价列表）
\end{itemize}

\textbf{2. 普通用户 --- 买家/卖家（USER）}

平台的主要使用者群体。同一用户可作为买家参与竞拍和购买，也可作为卖家发布商品。用户通过学号/工号注册后须经管理员审核身份。信誉分≥80的老用户发布商品免审核。

\begin{itemize}
    \item \textbf{作为买家}：浏览商品、搜索筛选、收藏商品、参与拍卖出价、与卖家IM私聊、下单购买、模拟支付、确认收货、评价交易
    \item \textbf{作为卖家}：发布商品、编辑商品信息、查看出价记录、与买家IM私聊、标记发货、查看订单、收到评价
    \item \textbf{管理个人信息}：编辑头像/昵称/手机号/简介、查看我的发布/我的竞拍/我的收藏/我的订单
\end{itemize}

\textbf{3. 管理员（ADMIN）}

平台运营维护者，负责审核用户身份、管理商品内容、监管拍卖行为、查看运营数据。

\begin{itemize}
    \item 用户管理：查看用户列表、审核新注册用户身份、封禁违规用户/解除封禁
    \item 商品管理：查看商品列表、审核通过新发布商品、下架违规商品
    \item 拍卖监管：查看进行中的拍卖列表、监控异常出价行为
    \item 分类管理：增删改查商品分类标签（树形结构，支持无限层级）
    \item 数据统计：仪表盘概览、交易趋势图（ECharts折线图）、分类占比（饼图）、用户活跃度（柱状图）、拍卖成交率
\end{itemize}

\subsubsection{用户对象关系图}

\begin{figure}[H]
\centering
{\small
\begin{verbatim}
                    ┌──────────────────────────┐
                    │    管理员 (ADMIN)          │
                    │   审核/监管/统计            │
                    └─────┬────────┬───────────┘
                审核身份    │        │   审核商品/监管拍卖
              ┌────────────┘        └──────────────┐
              ▼                                    ▼
  ┌───────────────────┐              ┌───────────────────────┐
  │   普通用户A        │              │   普通用户C             │
  │   (买家身份)       │◄──IM私聊───►│   (兼具买卖家身份)      │
  │                   │  参与竞拍    │                       │
  └───────────────────┘              └───────────────────────┘
                                        │         ▲
                                        │ 出售    │ 购买
                                        ▼         │
                              ┌───────────────────┐
                              │   普通用户B        │
                              │   (卖家身份)       │
                              └───────────────────┘

  ┌───────────┐    注册后成为    ┌───────────────────────┐
  │ 未登录访客  │ ──────────────► │  普通用户 (USER)       │
  │ ANONYMOUS  │                 └───────────────────────┘
  └───────────┘
\end{verbatim}
}
\caption{用户对象关系图}
\label{fig:user-relationship}
\end{figure}

\textbf{补充说明：}

\begin{itemize}
    \item \textbf{注册升级}：未登录访客通过注册并经过管理员审核后，升级为USER角色
    \item \textbf{一人双身份}：同一个USER账号可以在不同场景下作为买家或卖家
    \item \textbf{管理员封闭}：ADMIN不开放注册，由系统在数据库中直接指定
    \item \textbf{买卖家对称}：所有USER之间可以自由进行买卖
    \item \textbf{访客降级}：未登录访客只能单向浏览，无法与任何用户产生交易或消息交互
\end{itemize}

\subsubsection{功能模块与用户的对应一览表}

\begin{center}
\renewcommand{\arraystretch}{1.3}
\small
\begin{tabular}{|p{4.5cm}|c|c|c|c|}
\hline
\rowcolor{tableheader}
\textbf{功能模块} & \textbf{未登录访客} & \textbf{普通用户(买家)} & \textbf{普通用户(卖家)} & \textbf{管理员} \\
\hline
M1 用户注册与登录 & ● & ● & ● & ● \\
\hline
M1 用户资料管理 & --- & ● & ● & --- \\
\hline
M1 用户审核与封禁 & --- & --- & --- & ● \\
\hline
M2 商品浏览与搜索 & ● & ● & ● & ● \\
\hline
M2 商品发布 & --- & --- & ● & --- \\
\hline
M2 商品收藏 & --- & ● & ● & --- \\
\hline
M3 拍卖出价 & --- & ● & --- & --- \\
\hline
M3 拍卖成交判定 & --- & --- & --- & ---（系统自动） \\
\hline
M3 拍卖监管 & --- & --- & --- & ● \\
\hline
M4 IM即时通讯 & --- & ● & ● & --- \\
\hline
M5 订单创建 & --- & ● & --- & ---（拍卖成交系统自动） \\
\hline
M5 订单支付/收货 & --- & ● & --- & --- \\
\hline
M5 订单发货 & --- & --- & ● & --- \\
\hline
M6 商品审核与下架 & --- & --- & --- & ● \\
\hline
M6 分类标签管理 & --- & --- & --- & ● \\
\hline
M6 数据统计看板 & --- & --- & --- & ● \\
\hline
M6 交易评价 & --- & ● & ● & --- \\
\hline
\end{tabular}
\end{center}

\textbf{图例}：● 有权限 \quad --- 无权限

\textbf{补充说明：}

\begin{itemize}
    \item 拍卖成交判定（M3）为系统定时任务自动执行，无需任何用户角色触发
    \item 普通用户的"买家"和"卖家"身份是同一账号在不同操作中的角色切换
    \item 管理员拥有查看所有功能模块的权限，但不参与具体交易
\end{itemize}

% #################### 第3章 外部接口需求 ####################
\section{外部接口需求}

\subsection{用户接口}

使用图、表的方式，明确以下内容：

\begin{itemize}
    \item 需要给用户输出的数据
    \item 需要用户输入的数据
\end{itemize}

最好以UI图的方式展现。如果分解到各个功能模块说明的情况，此处仅需要明确指向：请参考...处说明。

\subsubsection{前端页面清单}

\begin{center}
\renewcommand{\arraystretch}{1.3}
\small
\begin{tabular}{|c|p{4cm}|p{6cm}|}
\hline
\rowcolor{tableheader}
\textbf{页面路由} & \textbf{页面名称} & \textbf{说明} \\
\hline
\verb|/| & 首页 & 商品Feed流 + 热门拍卖推荐 + 分类入口 \\
\hline
\verb|/login| & 登录页 & 学号 + 密码登录 \\
\hline
\verb|/register| & 注册页 & 学号 + 邮箱 + 密码 + 手机号 \\
\hline
\verb|/product/:id| & 商品详情页 & 图片轮播 + 一口价购买/拍卖竞价 \\
\hline
\verb|/product/create| & 发布商品页 & 表单：标题/描述/图片/分类/类型/价格/时长 \\
\hline
\verb|/search| & 搜索结果页 & 筛选栏 + 商品卡片列表 \\
\hline
\verb|/chat| & 消息列表页 & 会话列表 \\
\hline
\verb|/chat/:userId| & 聊天窗口 & 消息气泡 + 输入框 + 图片发送 \\
\hline
\verb|/user/:id| & 用户主页 & 信誉评分 + 交易记录 + 评价列表 \\
\hline
\verb|/user/profile| & 个人资料 & 编辑头像/昵称/手机号/简介 \\
\hline
\verb|/orders| & 我的订单 & 订单列表 + 状态筛选 \\
\hline
\verb|/admin| & 管理后台首页 & 仪表盘概览 \\
\hline
\verb|/admin/users| & 用户管理 & 审核列表 + 封禁操作 \\
\hline
\verb|/admin/products| & 商品管理 & 审核列表 + 下架操作 \\
\hline
\verb|/admin/stats| & 数据统计 & 图表看板 \\
\hline
\end{tabular}
\end{center}

\subsubsection{画面与功能对应详细说明}

\textbf{首页（P01）--- \texttt{/}}

显示商品Feed流（按发布时间倒序）、热门拍卖推荐（按浏览量排序）、分类快捷入口、顶部搜索栏。面向所有用户（含未登录）。特别要求：首次加载显示骨架屏，图片懒加载。

\textbf{登录页（P02）--- \texttt{/login}}

学号+密码登录表单。登录成功后跳转至首页。登录失败显示具体错误提示。特别要求：连续5次登录失败锁定账号15分钟。

\textbf{注册页（P03）--- \texttt{/register}}

学号、邮箱、密码（确认密码）、手机号注册表单。前端实时校验：学号为数字格式、邮箱格式、两次密码一致、密码长度≥8位含字母数字。提交后提示"注册成功，请等待管理员审核"。

\textbf{商品详情页（P04）--- \texttt{/product/:id}}

商品信息区域：图片轮播（点击放大）、标题、描述、分类标签、发布时间、浏览量、卖家信息卡片。交易区域分两种模式：一口价模式显示"立即购买"按钮；拍卖模式显示当前最高价（实时WebSocket更新）、出价输入框、拍卖倒计时、最新出价记录列表、"联系卖家"按钮。特别要求：拍卖模式下倒计时归零后，出价输入框自动隐藏。

\textbf{发布商品页（P05）--- \texttt{/product/create}}

表单：标题、描述（富文本）、图片上传（多图拖拽或点击上传，每张≤5MB）、分类选择（下拉树形）、交易模式切换（一口价/拍卖）、价格设置。特别要求：表单内容在localStorage自动暂存草稿。

\textbf{搜索结果页（P06）--- \texttt{/search}}

筛选面板（左侧）：分类树形筛选、价格区间、交易类型。排序栏（顶部）：最新发布/价格升序/价格降序/最多浏览。结果列表（右侧）：商品卡片分页展示。特别要求：搜索结果页URL可分享。

\textbf{IM消息列表页（P07）--- \texttt{/chat}}

会话列表（按最后消息时间倒序），每项显示对方头像+昵称+最后消息摘要+未读消息红点数字。空状态显示"暂无消息"插画。

\textbf{IM聊天窗口（P08）--- \texttt{/chat/:userId}}

消息气泡列表（己方靠右绿底，对方靠左白底）、文字输入框（Enter发送，Shift+Enter换行）、图片发送按钮、顶部显示对方信息+关联商品卡片。特别要求：WebSocket断开时消息通过HTTP降级发送，前端显示重连提示。

\textbf{用户主页（P09）--- \texttt{/user/:id}}

用户信息区域：头像、昵称、信誉分数值。交易统计：成交数量、好评率百分比。历史评价列表。

\textbf{个人资料页（P10）--- \texttt{/user/profile}}

头像上传、昵称、手机号、个人简介。保存按钮提交后更新。

\textbf{我的订单页（P11）--- \texttt{/orders}}

状态筛选Tab：全部 / 待付款 / 已发货 / 已完成。订单列表每项显示商品缩略图+标题+成交金额+状态标签+操作按钮。

\textbf{管理后台 --- \texttt{/admin/*}}（共3个管理页面，仅ADMIN角色可访问）

\begin{itemize}
    \item \texttt{/admin/users} 用户管理：用户列表（搜索+分页+状态筛选）、待审核用户高亮、"通过"/"驳回"按钮、封禁/解封操作
    \item \texttt{/admin/products} 商品管理：商品列表、"通过"/"下架"按钮
    \item \texttt{/admin/stats} 数据统计：ECharts图表---交易额趋势折线图、分类占比饼图、用户活跃度柱状图、拍卖成交率仪表盘
\end{itemize}

\subsubsection{交互设计原则与安全要求}

\textbf{交互设计原则：}
\begin{itemize}
    \item 响应式设计，适配桌面和移动端浏览器
    \item 拍卖页面关键操作（出价）按钮防重复点击
    \item 表单提交实时校验（前端 + 后端双重验证）
    \item 关键操作（发布商品/出价/下单）有确认提示
    \item WebSocket断线自动重连（指数退避）
\end{itemize}

\textbf{用户接口安全设计要求：}
\begin{itemize}
    \item 登录状态保持：JWT AccessToken（30分钟）存储于localStorage；RefreshToken（7天）存储于httpOnly Cookie
    \item 敏感操作确认：出价、下单支付、确认收货须弹出二次确认对话框
    \item 防越权：前端路由守卫 + 后端\verb|@PreAuthorize|双重校验
    \item 表单输入校验：前端Element Plus表单验证 + 后端Spring Validation双重校验
\end{itemize}

\subsection{硬件接口}

使用图、表的方式，明确以下内容：

\begin{itemize}
    \item 描述系统中软件和硬件每一接口的特征
\end{itemize}

本系统为纯Web应用，无特殊硬件接口需求。系统通过标准TCP/IP网络协议与客户端通信，运行于通用x86\_64架构服务器。以下硬件接口均不涉及：

\begin{itemize}
    \item 无串口（RS-232/RS-485）设备通信
    \item 无USB外设直接通信
    \item 无蓝牙/NFC外设通信
    \item 无PCI/PCIe设备驱动调用
    \item 无打印机、扫描仪等外接设备
\end{itemize}

输入设备为标准键盘和鼠标（桌面端）或触摸屏（移动端），输出设备为标准显示器或手机屏幕。

\subsection{软件接口}

使用图、表的方式，明确以下内容：

\begin{itemize}
    \item 描述该产品与其它外部组件的连接
    \item 明确并描述在软件组件之间交换数据或消息的目的
    \item 确定将在组件之间共享的数据
\end{itemize}

\subsubsection{数据库接口 --- MySQL 8.0}

\begin{itemize}
    \item \textbf{接口功能}：提供全部业务数据的持久化存储与查询
    \item \textbf{DBMS}：MySQL 8.0.36
    \item \textbf{连接方式}：JDBC，驱动程序 \verb|mysql-connector-j| 8.3.0
    \item \textbf{连接池}：HikariCP 5.1.0，最大连接数20，最小空闲5
    \item \textbf{连接参数}：\verb|jdbc:mysql://localhost:3306/campus_auction?useUnicode=true&characterEncoding=utf8mb4&serverTimezone=Asia/Shanghai|
    \item \textbf{ORM框架}：MyBatis-Plus 3.5.7（Lambda查询 + 分页插件）
    \item \textbf{迁移工具}：Flyway 10.x
\end{itemize}

\subsubsection{缓存接口 --- Redis 7}

\begin{itemize}
    \item \textbf{接口功能}：提供高性能缓存服务。具体场景：拍卖当前最高出价缓存、用户Session共享、在线用户状态维护、热门商品排行榜
    \item \textbf{Redis版本}：Redis 7.2
    \item \textbf{连接方式}：Redis Serialization Protocol (RESP)，客户端Lettuce 6.3
    \item \textbf{连接参数}：\verb|redis://localhost:6379|
    \item \textbf{数据结构使用}：String --- 拍卖当前价、拍卖得主；Set --- 在线用户集合；Sorted Set --- 热门商品排行
\end{itemize}

\subsubsection{文件存储接口}

\begin{itemize}
    \item \textbf{接口功能}：存储用户上传的商品图片
    \item \textbf{存储路径}：开发环境 \verb|./uploads/images/|，生产环境 \verb|/data/caip/uploads/images/|
    \item \textbf{文件命名规则}：\verb|{UUID}.{原始扩展名}|，防遍历
    \item \textbf{访问方式}：Spring Boot静态资源配置 \verb|/uploads/**| 映射到文件系统路径
    \item \textbf{预留接口}：\verb|FileStorageService| 接口可切换实现为阿里云OSS或MinIO对象存储
\end{itemize}

生产环境使用MinIO自建对象存储，或接入阿里云OSS。

\subsubsection{外部系统依赖}

本系统为完全独立部署，无外部系统依赖。

\begin{itemize}
    \item 支付功能：使用模拟支付，不接入微信支付/支付宝等第三方支付网关
    \item 短信/邮件：使用站内通知（notifications表 + WebSocket推送）代替
    \item 校园认证：学号+管理员人工审核方案，不接入学校统一身份认证系统
    \item 推荐算法：使用简单"热门排行"，不接入推荐引擎
\end{itemize}

\subsection{通讯接口}

使用图、表的方式，明确以下内容：

\begin{itemize}
    \item 描述与产品所使用的通信功能相关的需求
    \item 定义了相关的消息格式
    \item 规定通信安全或加密问题、数据传输速率和同步通信机制
\end{itemize}

\subsubsection{通讯接口总览}

\begin{center}
\renewcommand{\arraystretch}{1.3}
\begin{tabular}{|c|c|c|p{7cm}|}
\hline
\rowcolor{tableheader}
\textbf{接口类型} & \textbf{协议} & \textbf{端点格式} & \textbf{说明} \\
\hline
REST API & HTTPS & \verb|/api/v1/**| & JSON请求/响应，SpringDoc自动生成Swagger文档 \\
\hline
WebSocket（拍卖） & WSS & \verb|/ws/auction/{productId}| & 拍卖出价实时推送，STOMP Topic订阅模式 \\
\hline
WebSocket（IM） & WSS & \verb|/ws/chat| & IM实时通信，STOMP点对点模式 \\
\hline
\end{tabular}
\end{center}

\subsubsection{REST API 通讯接口}

\begin{itemize}
    \item \textbf{接口名称}：RESTful HTTP API
    \item \textbf{接口功能}：提供前端所有业务操作的HTTP接口
    \item \textbf{通信协议}：HTTPS (TLS 1.3)
    \item \textbf{数据格式}：请求/响应 JSON
    \item \textbf{认证方式}：JWT Bearer Token，请求头 \verb|Authorization: Bearer eyJhbG...|
    \item \textbf{API文档}：SpringDoc OpenAPI 2.6自动生成，访问 \verb|/swagger-ui.html|
\end{itemize}

\textbf{统一响应格式：}

\begin{lstlisting}[language=json]
{
  "code": 200,
  "message": "操作成功",
  "data": {},
  "timestamp": 1717507200000
}
\end{lstlisting}

\textbf{统一分页格式：}

\begin{lstlisting}[language=json]
{
  "total": 150,
  "page": 1,
  "size": 20,
  "items": []
}
\end{lstlisting}

\textbf{请求示例（登录）：}

\begin{lstlisting}
POST /api/v1/auth/login HTTP/1.1
Host: campus-auction.example.com
Content-Type: application/json

{
  "studentId": "2021300001",
  "password": "abc123"
}
\end{lstlisting}

\textbf{响应示例（登录成功）：}

\begin{lstlisting}[language=json]
{
  "code": 200,
  "message": "登录成功",
  "data": {
    "accessToken": "eyJhbGciOiJIUzI1NiJ9...",
    "refreshToken": "eyJhbGciOiJIUzI1NiJ9...",
    "userInfo": {
      "id": 1,
      "studentId": "2021300001",
      "username": "张三",
      "avatar": "/uploads/images/abc.jpg",
      "reputationScore": 100
    }
  },
  "timestamp": 1717507200000
}
\end{lstlisting}

\subsubsection{拍卖出价 WebSocket 接口}

\begin{itemize}
    \item \textbf{接口名称}：拍卖出价实时推送接口
    \item \textbf{接口功能}：实时推送某商品的出价更新给所有正在围观该商品的用户
    \item \textbf{通信协议}：WSS (WebSocket Secure) + STOMP子协议
    \item \textbf{通信方式}：发布-订阅，服务端发布到 \verb|/topic/auction/{productId}|，订阅者实时收到推送
    \item \textbf{连接认证}：STOMP CONNECT帧携带 \verb|Authorization: Bearer <token>| 头
\end{itemize}

\textbf{新出价推送消息格式：}

\begin{lstlisting}[language=json]
{
  "type": "NEW_BID",
  "payload": {
    "productId": 123,
    "bidId": 456,
    "bidAmount": 85.00,
    "currentPrice": 85.00,
    "bidderId": 10,
    "bidderName": "张三",
    "bidTime": "2026-06-09T14:30:00"
  },
  "timestamp": 1717933800000
}
\end{lstlisting}

\textbf{拍卖结束推送消息格式：}

\begin{lstlisting}[language=json]
{
  "type": "AUCTION_END",
  "payload": {
    "productId": 123,
    "winnerId": 10,
    "winnerName": "张三",
    "finalPrice": 85.00,
    "totalBids": 12
  },
  "timestamp": 1717935000000
}
\end{lstlisting}

\subsubsection{IM即时通讯 WebSocket 接口}

\begin{itemize}
    \item \textbf{接口名称}：IM即时通讯接口
    \item \textbf{接口功能}：实现买卖双方一对一实时私聊
    \item \textbf{通信协议}：WSS + STOMP
    \item \textbf{通信方式}：点对点，发送方发送到 \verb|/app/chat.send|，服务端转发到 \verb|/user/{userId}/queue/chat|
    \item \textbf{连接认证}：STOMP CONNECT帧携带Token
\end{itemize}

\textbf{IM消息格式：}

\begin{lstlisting}[language=json]
{
  "type": "NEW_MESSAGE",
  "payload": {
    "messageId": 789,
    "conversationId": 50,
    "senderId": 10,
    "senderName": "张三",
    "senderAvatar": "/uploads/images/abc.jpg",
    "content": "这个还在吗？",
    "contentType": "TEXT",
    "createdAt": "2026-06-09T14:35:00"
  },
  "timestamp": 1717934100000
}
\end{lstlisting}

\textbf{WebSocket统一消息信封：}

\begin{lstlisting}[language=json]
{
  "type": "NEW_BID | AUCTION_END | NEW_MESSAGE | NOTIFICATION",
  "payload": { },
  "timestamp": 1717507200000,
  "sender": "userId"
}
\end{lstlisting}

\subsubsection{CORS安全配置}

\begin{itemize}
    \item 允许来源：\verb|https://campus-auction.example.com|, \verb|http://localhost:5173|
    \item 允许方法：GET, POST, PUT, DELETE, OPTIONS
    \item 允许头部：Authorization, Content-Type
    \item 允许凭证：true
    \item 预检缓存：3600秒
\end{itemize}

% #################### 第4章 实现方式 ####################
\section{实现方式}

明确每个功能的实现方式，按照业务功能划分进行分类说明，有必要的情况分不同的章节进行说明。

\subsection{用户与认证功能 (M1)}

\subsubsection{功能概要}

文字描述用例概要，同时使用use-case图明确业务使用概要。

提供用户注册、登录、JWT无状态认证、个人信息管理功能。支持学号+邮箱注册，密码使用BCrypt加密存储（强度10）。通过Spring Security框架 + JWT实现基于角色的访问控制（RBAC）。

\textbf{认证流程：}

\begin{itemize}
    \item \textbf{注册流程}：用户提交注册信息 → 校验学号+邮箱唯一性 → BCryptPasswordEncoder加密密码 → 写入users表 → 返回注册成功
    \item \textbf{登录流程}：用户提交学号+密码 → UserService查询用户 → BCrypt.matches()校验密码 → JwtTokenProvider.generateToken()生成AccessToken(30分钟)和RefreshToken(7天) → 返回Token对
    \item \textbf{鉴权流程}：每个请求Authorization头携带Bearer Token → JwtAuthenticationFilter拦截 → 解析验证Token → 将用户信息写入SecurityContext → @PreAuthorize注解校验角色权限
\end{itemize}

\subsubsection{功能详细}

用文字或者图的方式对业务流程进行详细描述。本章要推荐使用的UML图是顺序图、活动图、状态图、协作图等。

\textbf{① 用户注册（M1-F01）}

\begin{itemize}
    \item 输入：RegisterDTO \{ studentId, email, password, phone \}
    \item 处理：校验学号和邮箱唯一性 → BCrypt加密密码 → 插入users表
    \item 输出：成功消息 + 提示"请等待管理员审核"
    \item 权限：无需登录
\end{itemize}

\textbf{② 用户登录（M1-F02）}

\begin{itemize}
    \item 输入：LoginDTO \{ studentId, password \}
    \item 处理：查询用户 → BCrypt校验密码 → 校验账号状态 → 生成JWT Token对 → 记录登录时间
    \item 输出：LoginVO \{ accessToken, refreshToken, userInfo \}
    \item 权限：无需登录
\end{itemize}

\textbf{③ Token刷新（M1-F03）}

\begin{itemize}
    \item 输入：RefreshToken字符串
    \item 处理：验证RefreshToken签名和有效期 → 提取userId → 生成新的AccessToken
    \item 输出：新AccessToken
    \item 权限：无需登录（但RefreshToken须有效）
\end{itemize}

\textbf{④ 获取个人信息（M1-F04）}

\begin{itemize}
    \item 输入：无（从JWT中提取userId）
    \item 处理：查询users表 → 组装UserVO
    \item 输出：UserVO
    \item 权限：须登录（USER或ADMIN）
\end{itemize}

\textbf{⑤ 编辑个人信息（M1-F05）}

\begin{itemize}
    \item 输入：UpdateUserDTO \{ nickname, phone, avatar \}
    \item 处理：校验手机号格式 → 如有新头像则上传文件返回URL → 更新users表
    \item 输出：更新后的UserVO
    \item 权限：须登录，只能编辑自己的信息
\end{itemize}

\textbf{⑥ 查看用户主页（M1-F06）}

\begin{itemize}
    \item 输入：userId（URL路径参数）
    \item 处理：查询用户公开信息 → 统计交易数量 → 统计好评率 → 查询最近10条评价
    \item 输出：UserProfileVO
    \item 权限：所有用户（含未登录）
\end{itemize}

\textbf{⑦ 管理员-身份审核（M1-F07）}

\begin{itemize}
    \item 输入：userId, 审核结果（APPROVED/REJECTED）
    \item 处理：@PreAuthorize校验 → UPDATE users SET status='ACTIVE' 或保持'PENDING' + 发送驳回通知
    \item 输出：成功消息 + WebSocket推送通知
    \item 权限：仅ADMIN
\end{itemize}

\textbf{⑧ 管理员-封禁/解封（M1-F08）}

\begin{itemize}
    \item 输入：userId, 操作类型（BAN/UNBAN）
    \item 处理：ADMIN权限校验 → UPDATE users SET status='BANNED'/'ACTIVE'
    \item 输出：成功消息
    \item 权限：仅ADMIN
\end{itemize}

\subsubsection{相关功能}

本功能如果需要多个功能才能实现或本功能和其他组件有关联的情况下，要明确此关联关系。

\begin{center}
\renewcommand{\arraystretch}{1.3}
\begin{tabular}{|p{3cm}|p{4cm}|p{7cm}|}
\hline
\rowcolor{tableheader}
\textbf{组件名} & \textbf{说明} & \textbf{关系} \\
\hline
M3 拍卖引擎 & 拍卖出价操作 & 所有出价操作须通过JwtAuthenticationFilter验证 \\
\hline
M4 IM通讯 & 即时通讯 & 发送消息时从SecurityContext获取当前用户ID \\
\hline
M5 订单交易 & 订单操作 & 支付/发货/收货操作须验证操作者为本人 \\
\hline
M6 管理后台 & 管理功能 & 所有/admin/**接口标注@PreAuthorize("hasRole('ADMIN')") \\
\hline
\end{tabular}
\end{center}

\subsubsection{用户接口}

\begin{center}
\renewcommand{\arraystretch}{1.3}
\begin{tabular}{|c|c|c|p{5cm}|}
\hline
\rowcolor{tableheader}
\textbf{接口} & \textbf{方法} & \textbf{认证} & \textbf{说明} \\
\hline
\verb|/api/v1/auth/register| & POST & 无 & 注册新用户 \\
\hline
\verb|/api/v1/auth/login| & POST & 无 & 登录，返回JWT Token对 \\
\hline
\verb|/api/v1/auth/refresh| & POST & 无 & 刷新AccessToken \\
\hline
\verb|/api/v1/users/me| & GET & 需 & 获取当前登录用户信息 \\
\hline
\verb|/api/v1/users/me| & PUT & 需 & 编辑个人信息 \\
\hline
\verb|/api/v1/users/{id}| & GET & 需 & 查看任意用户公开主页 \\
\hline
\end{tabular}
\end{center}

\begin{figure}[H]
\centering
{\small
\begin{verbatim}
┌─────────────────────────────────────────────────────────┐
│                    用户模块类图                           │
│                                                         │
│  ┌──────────────────┐    ┌─────────────────────────┐    │
│  │  AuthController   │    │     UserController       │    │
│  │  +register()      │    │     +getMyProfile()      │    │
│  │  +login()         │    │     +updateMyProfile()   │    │
│  │  +refreshToken()  │    │     +getUserProfile()    │    │
│  └────────┬─────────┘    └───────────┬─────────────┘    │
│           │                          │                  │
│           └──────────┬───────────────┘                  │
│                      ▼                                  │
│           ┌──────────────────────┐                      │
│           │   UserServiceImpl     │                      │
│           │   -userMapper         │                      │
│           │   -passwordEncoder    │                      │
│           │   +register()         │                      │
│           │   +authenticate()     │                      │
│           │   +findById()         │                      │
│           └──────────┬───────────┘                      │
│                      │                                  │
│           ┌──────────┴───────────┐                      │
│           ▼                      ▼                      │
│  ┌────────────────┐   ┌────────────────────┐           │
│  │ JwtTokenProvider│   │   SecurityConfig    │           │
│  │ +generateToken()│   │   +filterChain()    │           │
│  │ +validateToken()│   │   +passwordEncoder()│           │
│  └────────────────┘   └────────────────────┘           │
└─────────────────────────────────────────────────────────┘
\end{verbatim}
}
\caption{用户模块类图}
\label{fig:m1-class}
\end{figure}

\begin{figure}[H]
\centering
{\small
\begin{verbatim}
  客户端(Vue)         AuthController      UserServiceImpl    JwtTokenProvider      MySQL
     │                     │                    │                  │                 │
     │ POST /login         │                    │                  │                 │
     │────────────────────►│                    │                  │                 │
     │                     │ authenticate()     │                  │                 │
     │                     │───────────────────►│                  │                 │
     │                     │                    │ SELECT * FROM    │                 │
     │                     │                    │ users WHERE      │                 │
     │                     │                    │ student_id=?     │                 │
     │                     │                    │─────────────────────────────────►│
     │                     │                    │◄─────────────────────────────────│
     │                     │                    │                  │                 │
     │                     │              [验证成功]                 │                 │
     │                     │                    │ generateAccessToken()            │
     │                     │                    │─────────────────►│                │
     │                     │                    │◄─────────────────│ accessToken    │
     │                     │                    │ generateRefreshToken()           │
     │                     │                    │─────────────────►│                │
     │                     │                    │◄─────────────────│ refreshToken   │
     │                     │◄───────────────────│                  │                 │
     │◄────────────────────│ LoginVO{token,userInfo}              │                 │
     │ 200 登录成功         │                    │                  │                 │
\end{verbatim}
}
\caption{登录认证时序图}
\label{fig:m1-sequence}
\end{figure}

\subsection{商品发布与搜索功能 (M2)}

\subsubsection{功能概要}

提供商品发布（一口价和拍卖两种模式）、商品列表浏览、全文搜索、多条件筛选、商品收藏功能。支持多图上传和分类标签管理。新用户发布商品需管理员审核，信誉分≥80的老用户自动审核通过。

\subsubsection{功能详细}

\textbf{① 发布商品（M2-F01）}

\begin{itemize}
    \item 输入：ProductCreateDTO \{ title, description, images[], categoryId, type(SALE/AUCTION), price, auctionDuration? \}
    \item 处理：校验图片类型和大小 → 逐个上传图片返回URL数组 → 创建Product记录 → 关联product\_images记录 → 信誉分≥80自动审核通过
    \item 输出：ProductVO
    \item 权限：USER
\end{itemize}

\textbf{② 商品列表（M2-F02）}

\begin{itemize}
    \item 输入：page, size, sort(created\_at/price/view\_count)
    \item 处理：分页查询status='ACTIVE'的商品 → 关联卖家信息和首张图片
    \item 输出：Page<ProductCardVO>
    \item 权限：无需认证
\end{itemize}

\textbf{③ 商品详情（M2-F03）}

\begin{itemize}
    \item 输入：productId（URL路径参数）
    \item 处理：查询product + product\_images + seller信息 → 若type=AUCTION则关联最新出价记录 → view\_count+1
    \item 输出：ProductDetailVO
    \item 权限：无需认证
\end{itemize}

\textbf{④ 全文搜索（M2-F04）}

\begin{itemize}
    \item 输入：q(关键词), categoryId, minPrice, maxPrice, type, sort, page, size
    \item 处理：MySQL ngram全文索引搜索title+description → 多条件动态SQL筛选 → 分页排序。若全文索引无结果，降级为LIKE模糊搜索
    \item 输出：Page<ProductCardVO>
    \item 权限：无需认证
\end{itemize}

\textbf{搜索SQL（核心查询逻辑）：}

\begin{lstlisting}[language=sql]
SELECT p.*, pi.image_url AS cover_image
FROM products p
LEFT JOIN product_images pi ON p.id = pi.product_id AND pi.sort_order = 0
WHERE p.status = 'ACTIVE'
  AND MATCH(p.title, p.description) AGAINST(? IN BOOLEAN MODE)
  AND (? IS NULL OR p.category_id = ?)
  AND (? IS NULL OR p.current_price BETWEEN ? AND ?)
  AND (? IS NULL OR p.type = ?)
ORDER BY CASE ?
  WHEN 'price_asc' THEN p.current_price END ASC,
  CASE ? WHEN 'created_at' THEN p.created_at END DESC
LIMIT ?, ?
\end{lstlisting}

MySQL全文索引配置：

\begin{lstlisting}[language=sql]
ALTER TABLE products ADD FULLTEXT INDEX ft_products_search
    (title, description) WITH PARSER ngram;
SET GLOBAL ngram_token_size = 2;
\end{lstlisting}

\textbf{⑤ 收藏/取消收藏（M2-F05）}

\begin{itemize}
    \item 输入：productId
    \item 处理：查询product\_favorites表是否已有记录 → 有则DELETE（取消），无则INSERT（收藏）
    \item 输出：当前收藏状态（true/false）
    \item 权限：USER
\end{itemize}

\textbf{⑥ 我的收藏（M2-F06）}

\begin{itemize}
    \item 输入：page, size
    \item 处理：分页查询当前用户收藏的商品
    \item 输出：Page<ProductCardVO>
    \item 权限：USER
\end{itemize}

\subsubsection{相关功能}

\begin{center}
\renewcommand{\arraystretch}{1.3}
\begin{tabular}{|p{3cm}|p{4cm}|p{7cm}|}
\hline
\rowcolor{tableheader}
\textbf{组件名} & \textbf{说明} & \textbf{关系} \\
\hline
M3 拍卖引擎 & 拍卖商品 & type='AUCTION'的商品发布后，M3模块读取auction\_start/auction\_end启动拍卖 \\
\hline
M5 订单交易 & 一口价购买 & type='SALE'商品被购买后调用OrderService.createOrder() \\
\hline
M6 管理后台 & 商品审核 & 管理员审核商品 \\
\hline
\end{tabular}
\end{center}

\subsubsection{用户接口}

\begin{center}
\renewcommand{\arraystretch}{1.3}
\begin{tabular}{|c|c|p{6cm}|}
\hline
\rowcolor{tableheader}
\textbf{接口} & \textbf{方法} & \textbf{说明} \\
\hline
\verb|/api/v1/products| & GET & 商品列表（分页+排序） \\
\hline
\verb|/api/v1/products| & POST & 发布商品 \\
\hline
\verb|/api/v1/products/{id}| & GET & 商品详情 \\
\hline
\verb|/api/v1/products/{id}| & PUT & 编辑商品（仅卖家本人） \\
\hline
\verb|/api/v1/products/search| & GET & 搜索筛选 \\
\hline
\verb|/api/v1/products/{id}/favorite| & POST & 收藏/取消 \\
\hline
\verb|/api/v1/favorites| & GET & 我的收藏 \\
\hline
\end{tabular}
\end{center}

% M2 class diagram figure
\begin{figure}[H]
\centering
{\small
\begin{verbatim}
┌───────────────────────────────────────────────────────────────┐
│                      商品模块类图                              │
│                                                               │
│  ┌─────────────────────┐                                      │
│  │  ProductController   │                                      │
│  │  +list() +create()   │                                      │
│  │  +getById()          │                                      │
│  │  +search()           │                                      │
│  │  +toggleFavorite()   │                                      │
│  └──────────┬──────────┘                                      │
│             │                                                  │
│             ▼                                                  │
│  ┌──────────────────────────┐                                 │
│  │   ProductServiceImpl      │                                 │
│  │   -productMapper          │                                 │
│  │   -imageMapper            │                                 │
│  │   -fileService            │                                 │
│  │   +create() +getDetail()  │                                 │
│  │   +search() +toggleFav()  │                                 │
│  └───────┬──────────────────┘                                 │
│          │                                                     │
│  ┌───────┴──────────┬──────────────────┐                      │
│  ▼                  ▼                  ▼                       │
│ ┌──────────┐  ┌──────────────┐  ┌─────────────────┐          │
│ │ Product   │  │ FileStorage   │  │ ProductMapper    │          │
│ │ -id       │  │ Service       │  │ +selectPage()    │          │
│ │ -sellerId │  │ <<interface>> │  │ +fulltextSearch()│          │
│ │ -title    │  │ +upload()     │  └─────────────────┘          │
│ │ -type     │  │ +delete()     │                               │
│ │ -price    │  └──────────────┘                               │
│ │ -status   │                                                  │
│ └──────────┘                                                  │
└───────────────────────────────────────────────────────────────┘
\end{verbatim}
}
\caption{商品模块类图}
\label{fig:m2-class}
\end{figure}

\subsection{实时拍卖引擎功能 (M3)}

\subsubsection{功能概要}

实现英式升价拍卖。买家对进行中的拍卖商品出价，系统通过Redis Lua脚本保证出价操作的原子性，再通过WebSocket (STOMP Topic)将出价实时推送给所有正在围观该商品的用户。拍卖截止后由@Scheduled定时任务自动判定得主。

\textbf{拍卖生命周期：}

卖家发布拍卖商品（type=AUCTION, auction\_start=NOW, auction\_end=NOW+时长）
→ 拍卖进行中（接受用户出价）
→ auction\_end到达
→ 有人出价：winner=最高出价者, status=SOLD → 生成订单
→ 无人出价：status=EXPIRED → 流拍

\subsubsection{Redis Lua出价脚本}

出价操作使用Redis Lua脚本保证原子性，一个Lua调用中完成：读取当前价 → 校验出价>当前价 → 更新当前价和最高出价者。

\begin{lstlisting}[language=lua]
-- auction_place_bid.lua
-- KEYS[1] = auction:{productId}:price
-- KEYS[2] = auction:{productId}:winner
-- ARGV[1] = bidAmount      ARGV[2] = userId
-- ARGV[3] = startPrice     ARGV[4] = sellerId
-- 返回: {0, 错误信息} 或 {1, 新价格}

local currentPrice = redis.call('GET', KEYS[1])
if currentPrice == false then
    currentPrice = ARGV[3]  -- 首次出价，使用起拍价
end

local bidAmount = tonumber(ARGV[1])

-- 卖家不能给自己的商品出价
if ARGV[2] == ARGV[4] then
    return {0, "不能给自己的商品出价"}
end

-- 出价必须高于当前价
if bidAmount <= tonumber(currentPrice) then
    return {0, "出价必须高于当前价" .. currentPrice}
end

redis.call('SET', KEYS[1], ARGV[1])
redis.call('SET', KEYS[2], ARGV[2])
return {1, ARGV[1]}
\end{lstlisting}

\subsubsection{功能详细}

\textbf{① 出价（M3-F02）---核心功能}

\begin{itemize}
    \item 输入：productId, bidAmount
    \item 处理：①校验商品status=ACTIVE且type=AUCTION ②校验拍卖未结束 ③校验出价者不是卖家本人 ④执行Redis Lua脚本原子校验+更新 ⑤成功 → 通过SimpMessagingTemplate广播给所有订阅者 ⑥@Async异步持久化bids表
    \item 输出：BidVO
    \item 权限：USER（须登录）
\end{itemize}

\textbf{② 实时推送（M3-F03）}

\begin{itemize}
    \item 通信方式：STOMP Topic \verb|/topic/auction/{productId}|
    \item 消息类型：NEW\_BID（新出价）、AUCTION\_END（拍卖结束/成交）、AUCTION\_EXPIRED（流拍）
    \item 订阅者：前端页面打开商品详情时自动订阅，离开页面时取消订阅
\end{itemize}

\textbf{③ 定时判定成交（M3-F05）}

\begin{itemize}
    \item 触发：@Scheduled(fixedRate=60000) 每分钟执行一次
    \item 处理：SELECT到期拍卖 → 从Redis读取最高价和winner → 若有winner→UPDATE products SET status='SOLD' → 调用OrderService.createOrder() → WebSocket广播AUCTION\_END；若无winner → UPDATE products SET status='EXPIRED' → 广播流拍
    \item 防重复执行：Redis SETNX分布式锁 \verb|auction:check:lock|，过期60秒
    \item 权限：系统定时任务
\end{itemize}

\textbf{④ 查看出价记录（M3-F04）}

\begin{itemize}
    \item 输入：productId, page, size
    \item 处理：分页查询bids表（ORDER BY bid\_time DESC）
    \item 输出：Page<BidVO>
    \item 权限：无需认证
\end{itemize}

\textbf{⑤ 获取当前最高价（M3-F06）}

\begin{itemize}
    \item 输入：productId
    \item 处理：redisTemplate.opsForValue().get("auction:\{productId\}:price")
    \item 输出：BigDecimal
    \item 权限：无需认证
\end{itemize}

\subsubsection{相关功能}

\begin{center}
\renewcommand{\arraystretch}{1.3}
\begin{tabular}{|p{3cm}|p{4cm}|p{7cm}|}
\hline
\rowcolor{tableheader}
\textbf{组件名} & \textbf{说明} & \textbf{关系} \\
\hline
M2 商品管理 & 拍卖商品基础信息 & 读取商品起拍价和时长；成交/流拍后更新product.status和winner\_id \\
\hline
M5 订单交易 & 拍卖成交后生成订单 & 成交后调用orderService.createOrder() \\
\hline
M1 用户认证 & 出价者身份 & 出价须JWT验证 \\
\hline
M4 IM通讯 & 通知买卖双方 & 拍卖成交/流拍后通知买卖双方 \\
\hline
\end{tabular}
\end{center}

\subsubsection{用户接口}

\begin{center}
\renewcommand{\arraystretch}{1.3}
\begin{tabular}{|c|c|p{6cm}|}
\hline
\rowcolor{tableheader}
\textbf{接口} & \textbf{方法} & \textbf{说明} \\
\hline
\verb|/api/v1/products/{id}/bids| & GET & 出价记录（分页） \\
\hline
\verb|/api/v1/products/{id}/bids| & POST & 提交出价 \\
\hline
\verb|/api/v1/products/{id}/current-price| & GET & 当前最高价 \\
\hline
\verb|/ws/auction/{productId}| & WS(订阅) & 实时推送（STOMP订阅） \\
\hline
\end{tabular}
\end{center}

\begin{figure}[H]
\centering
{\small
\begin{verbatim}
  买家浏览器     AuctionController  AuctionServiceImpl    Redis     STOMP Broker  其他用户
     │                │                   │                │           │           │
     │ POST /bids     │                   │                │           │           │
     │───────────────►│                   │                │           │           │
     │                │  placeBid()       │                │           │           │
     │                │──────────────────►│                │           │           │
     │                │                   │ 校验商品状态    │           │           │
     │                │                   │ EVAL lua脚本    │           │           │
     │                │                   │───────────────►│           │           │
     │                │                   │◄───────────────│ {1,85.00} │           │
     │                │                   │ convertAndSend()            │           │
     │                │                   │────────────────────────────►│           │
     │                │                   │                │           │ NEW_BID   │
     │                │                   │                │           │──────────►│
     │                │                   │ @Async persist  │           │           │
     │                │                   │ INSERT bids     │           │           │
     │                │◄──────────────────│                 │           │           │
     │◄───────────────│ BidVO{id,amount}  │                 │           │           │
     │ 200 出价成功    │                   │                 │           │           │
\end{verbatim}
}
\caption{拍卖出价时序图}
\label{fig:m3-bid-sequence}
\end{figure}

\begin{figure}[H]
\centering
{\small
\begin{verbatim}
  @Scheduled(每分钟)  AuctionServiceImpl     MySQL        Redis    OrderService  WebSocket
     │                    │                   │            │          │            │
     │ checkExpired()     │                   │            │          │            │
     │───────────────────►│                   │            │          │            │
     │                    │ SELECT 到期拍卖    │            │          │            │
     │                    │──────────────────►│            │          │            │
     │                    │◄──────────────────│            │          │            │
     │                    │                   │            │          │            │
     │                    │ GET auction:{id}:winner        │          │            │
     │                    │───────────────────────────────►│          │            │
     │              [winnerId != null: 成交]               │          │            │
     │                    │ UPDATE products SET status='SOLD'          │            │
     │                    │──────────────────►│            │          │            │
     │                    │ createOrder()     │            │          │            │
     │                    │───────────────────────────────────────────►│            │
     │                    │                   │            │ INSERT orders           │
     │                    │ 广播 AUCTION_END  │            │          │            │
     │                    │─────────────────────────────────────────────────────►│
     │              [winnerId == null: 流拍]               │          │            │
     │                    │ UPDATE products SET status='EXPIRED'       │            │
     │                    │──────────────────►│            │          │            │
     │                    │ 广播 AUCTION_EXPIRED           │          │            │
     │                    │─────────────────────────────────────────────────────►│
     │                    │ DEL auction:{id}:price winner  │          │            │
     │                    │───────────────────────────────►│          │            │
\end{verbatim}
}
\caption{拍卖成交判定时序图}
\label{fig:m3-end-sequence}
\end{figure}

\subsection{即时通讯功能 (M4)}

\subsubsection{功能概要}

实现买卖双方一对一文字+图片私聊功能。使用WebSocket(STOMP)协议进行实时消息收发，同时提供HTTP REST接口作为WebSocket不可用时的降级方案。聊天记录持久化存储在MySQL中。消息已读/未读状态实时更新。

\subsubsection{功能详细}

\textbf{① 获取会话列表（M4-F01）}

\begin{itemize}
    \item 输入：无（当前用户从SecurityContext获取）
    \item 处理：查询当前用户参与的所有conversations → 关联对方用户信息 → 计算未读消息数 → 按last\_message\_at倒序
    \item 输出：List<ConversationVO>（每项含对方头像/昵称、最后消息摘要、未读计数）
    \item 权限：USER
\end{itemize}

\textbf{② 创建会话（M4-F02）}

\begin{itemize}
    \item 输入：targetUserId, productId（可选）
    \item 处理：检查是否已存在相同会话 → 无则创建(user1\_id=MIN, user2\_id=MAX保证唯一)
    \item 输出：ConversationVO
    \item 权限：USER
\end{itemize}

\textbf{③ 获取消息历史（M4-F03）}

\begin{itemize}
    \item 输入：conversationId, page, size
    \item 处理：分页查询messages表（ORDER BY created\_at ASC）→ 标记对方消息为已读
    \item 输出：Page<MessageVO>
    \item 权限：USER（须是会话参与方）
\end{itemize}

\textbf{④ 发送消息 --- WebSocket（M4-F04，主要方式）}

\begin{itemize}
    \item 输入：STOMP SEND to \verb|/app/chat.send| \{ conversationId, content, contentType \}
    \item 处理：从SecurityContext获取senderId → 校验sender属于会话参与方 → messageMapper.insert()持久化MySQL → messagingTemplate.convertAndSendToUser()推送给接收方 → 回执给发送方
    \item 输出：实时WebSocket推送
    \item 权限：USER
\end{itemize}

\textbf{⑤ 发送消息 --- HTTP降级（M4-F05，备用方式）}

\begin{itemize}
    \item 输入：POST \{ conversationId, content, contentType \}
    \item 处理：同上，WebSocket不可用时的HTTP降级方案
    \item 输出：MessageVO
    \item 权限：USER
\end{itemize}

\subsubsection{相关功能}

\begin{center}
\renewcommand{\arraystretch}{1.3}
\begin{tabular}{|p{3cm}|p{4cm}|p{7cm}|}
\hline
\rowcolor{tableheader}
\textbf{组件名} & \textbf{说明} & \textbf{关系} \\
\hline
M1 用户认证 & 消息发送者身份 & 发送消息时从SecurityContext获取sender\_id \\
\hline
M2 商品浏览 & 从商品发起对话 & 从商品详情页"联系卖家"按钮创建带product\_id的会话 \\
\hline
M3 拍卖引擎 & 系统通知 & 拍卖成交/流拍后自动发送系统通知消息 \\
\hline
\end{tabular}
\end{center}

\subsubsection{用户接口}

\begin{center}
\renewcommand{\arraystretch}{1.3}
\begin{tabular}{|c|c|p{6cm}|}
\hline
\rowcolor{tableheader}
\textbf{接口} & \textbf{方法} & \textbf{说明} \\
\hline
\verb|/api/v1/conversations| & GET & 会话列表（含未读数） \\
\hline
\verb|/api/v1/conversations| & POST & 创建会话 \\
\hline
\verb|/api/v1/conversations/{id}/messages| & GET & 消息历史（分页，正序） \\
\hline
\verb|/api/v1/conversations/{id}/messages| & POST & 发送消息（HTTP降级） \\
\hline
\verb|/api/v1/conversations/{id}/read| & PUT & 标记已读 \\
\hline
\verb|/ws/chat| & WS(双向) & IM实时通信（STOMP双向，主要方式） \\
\hline
\end{tabular}
\end{center}

\begin{figure}[H]
\centering
{\small
\begin{verbatim}
  发送方浏览器      STOMP Broker    ChatController   ChatServiceImpl    MySQL   接收方浏览器
     │                  │               │                │              │          │
     │ STOMP SEND       │               │                │              │          │
     │ /app/chat.send   │               │                │              │          │
     │─────────────────►│               │                │              │          │
     │                  │ @MessageMapping                │              │          │
     │                  │──────────────►│                │              │          │
     │                  │               │ sendMessage()  │              │          │
     │                  │               │───────────────►│              │          │
     │                  │               │                │ INSERT msg   │          │
     │                  │               │                │─────────────►│          │
     │                  │               │                │◄─────────────│ msgId=789│
     │                  │               │                │              │          │
     │                  │               │ convertAndSendToUser("15")   │          │
     │                  │◄──────────────│◄──────────────│              │          │
     │                  │ STOMP MESSAGE │                │              │          │
     │                  │─────────────────────────────────────────────────────────►│
     │                  │ {type:"NEW_MESSAGE", payload:{msgId:789,...}}            │
     │                  │                │                │              │          │
     │                  │ 发送回执       │                │              │          │
     │◄─────────────────│ {type:"MESSAGE_SENT", payload:{msgId:789}}              │
\end{verbatim}
}
\caption{IM消息收发时序图}
\label{fig:m4-sequence}
\end{figure}

\subsection{订单与交易功能 (M5)}

\subsubsection{功能概要}

管理订单全生命周期。订单来源有两类：(a)拍卖成交后由M3模块自动生成；(b)一口价商品由买家点击"立即购买"后生成。实现订单状态机，严格控制各状态间的合法跳转。提供模拟支付功能。

\textbf{订单状态机：}

\begin{verbatim}
[下单] → PENDING(待付款)
              ├──→ PAID(已付款) [买家模拟支付]
              └──→ CANCELLED(已取消) [买家取消，仅PENDING状态]
         PAID → SHIPPED(已发货) [卖家标记发货]
         SHIPPED → COMPLETED(已完成) [买家确认收货] → 可评价
\end{verbatim}

\textbf{状态跳转校验规则：}

\begin{itemize}
    \item PENDING → PAID：允许，操作者=买家
    \item PENDING → CANCELLED：允许，操作者=买家
    \item PAID → SHIPPED：允许，操作者=卖家
    \item SHIPPED → COMPLETED：允许，操作者=买家
    \item 其他跳转：禁止
\end{itemize}

\begin{figure}[H]
\centering
{\small
\begin{verbatim}
                    订单状态机图

     [*] ──→ PENDING ──→ PAID ──→ SHIPPED ──→ COMPLETED ──→ [*]
               │                                      (可评价)
               │
               └──→ CANCELLED ──→ [*]
               (买家取消)
\end{verbatim}
}
\caption{订单状态机图}
\label{fig:m5-state}
\end{figure}

\subsubsection{功能详细}

\textbf{① 创建订单（M5-F01）}

\begin{itemize}
    \item 输入：productId（一口价购买时由买家触发；拍卖成交时由M3调用）
    \item 处理：查询商品信息 → 创建Order（status=PENDING）→ 更新商品status='SOLD'
    \item 输出：OrderVO
    \item 权限：USER(买家) 或 系统（M3调用）
\end{itemize}

\textbf{② 模拟支付（M5-F02）}

\begin{itemize}
    \item 输入：orderId
    \item 处理：校验status=PENDING → 校验操作者为买家本人 → UPDATE status='PAID', paid\_at=NOW()
    \item 输出：OrderVO
    \item 权限：USER(买家本人)
\end{itemize}

\textbf{③ 标记发货（M5-F03）}

\begin{itemize}
    \item 输入：orderId
    \item 处理：校验status=PAID → 校验操作者为卖家本人 → UPDATE status='SHIPPED'
    \item 输出：OrderVO
    \item 权限：USER(卖家本人)
\end{itemize}

\textbf{④ 确认收货（M5-F04）}

\begin{itemize}
    \item 输入：orderId
    \item 处理：校验status=SHIPPED → 校验操作者为买家本人 → UPDATE status='COMPLETED'
    \item 输出：OrderVO + 触发评价入口
    \item 权限：USER(买家本人)
\end{itemize}

\textbf{⑤ 取消订单（M5-F05）}

\begin{itemize}
    \item 输入：orderId
    \item 处理：校验status=PENDING → 校验操作者为买家本人 → UPDATE status='CANCELLED' → UPDATE products SET status='ACTIVE'
    \item 输出：OrderVO
    \item 权限：USER(买家本人)
\end{itemize}

\textbf{⑥ 订单列表（M5-F06）}

\begin{itemize}
    \item 输入：状态筛选, page, size
    \item 处理：按角色查询 → 关联商品信息
    \item 输出：Page<OrderVO>
    \item 权限：USER
\end{itemize}

\subsubsection{相关功能}

\begin{center}
\renewcommand{\arraystretch}{1.3}
\begin{tabular}{|p{3cm}|p{4cm}|p{7cm}|}
\hline
\rowcolor{tableheader}
\textbf{组件名} & \textbf{说明} & \textbf{关系} \\
\hline
M2 商品管理 & 商品信息 & 读取商品信息；更新商品状态 \\
\hline
M3 拍卖引擎 & 拍卖成交 & 拍卖成交后M3调用createOrder() \\
\hline
M6 信誉评价 & 评价系统 & 订单COMPLETED后双方可评价 \\
\hline
\end{tabular}
\end{center}

\subsubsection{用户接口}

\begin{center}
\renewcommand{\arraystretch}{1.3}
\begin{tabular}{|c|c|p{6cm}|}
\hline
\rowcolor{tableheader}
\textbf{接口} & \textbf{方法} & \textbf{说明} \\
\hline
\verb|/api/v1/orders| & GET & 订单列表 \\
\hline
\verb|/api/v1/orders| & POST & 创建订单（一口价购买） \\
\hline
\verb|/api/v1/orders/{id}| & GET & 订单详情（仅参与方可查看） \\
\hline
\verb|/api/v1/orders/{id}/pay| & POST & 模拟支付（仅买家） \\
\hline
\verb|/api/v1/orders/{id}/ship| & POST & 标记发货（仅卖家） \\
\hline
\verb|/api/v1/orders/{id}/confirm| & POST & 确认收货（仅买家） \\
\hline
\verb|/api/v1/orders/{id}/cancel| & POST & 取消订单（仅买家，仅PENDING状态） \\
\hline
\end{tabular}
\end{center}

\begin{figure}[H]
\centering
{\small
\begin{verbatim}
┌───────────────────────────────────────────────────────────────┐
│                      订单模块类图                              │
│                                                               │
│  ┌──────────────────────────┐                                 │
│  │    OrderController        │                                 │
│  │    +createOrder()         │                                 │
│  │    +listOrders()          │                                 │
│  │    +pay() +ship()         │                                 │
│  │    +confirm() +cancel()   │                                 │
│  └───────────┬──────────────┘                                 │
│              │                                                 │
│              ▼                                                 │
│  ┌────────────────────────────────┐                           │
│  │     OrderServiceImpl            │                           │
│  │     -orderMapper                │                           │
│  │     +createOrder() +pay()       │                           │
│  │     +ship() +confirm()          │                           │
│  │     -validateTransition()       │                           │
│  └────────┬───────────────────────┘                           │
│           │                                                    │
│  ┌────────┴──────────┐    ┌────────────────────┐              │
│  │ OrderStatusValidator│    │ OrderStatusEnum    │              │
│  │ +validate()        │    │ PENDING, PAID,      │              │
│  │ -ALLOWED_TRANSITIONS│   │ SHIPPED, COMPLETED, │              │
│  └────────────────────┘    │ CANCELLED           │              │
│                             └────────────────────┘              │
└───────────────────────────────────────────────────────────────┘
\end{verbatim}
}
\caption{订单模块类图}
\label{fig:m5-class}
\end{figure}

\subsection{管理后台与统计功能 (M6)}

\subsubsection{功能概要}

提供管理员专用后台功能：用户身份审核与封禁管理、商品审核与下架管理、分类标签管理、拍卖异常监控、平台数据统计看板（ECharts可视化）。所有管理接口统一标注 \verb|@PreAuthorize("hasRole('ADMIN')")|。

\subsubsection{功能详细}

\textbf{① 用户列表与审核（M6-F01/F02）}

\begin{itemize}
    \item 输入：page, size, status筛选, keyword搜索
    \item 处理：分页查询users表 → 支持按status和keyword筛选 → 审核通过：UPDATE status='ACTIVE' + 通知用户
    \item 输出：Page<UserVO>
    \item 权限：ADMIN
\end{itemize}

\textbf{② 封禁/解封用户（M6-F03）}

\begin{itemize}
    \item 输入：userId, 操作类型(BAN/UNBAN)
    \item 处理：封禁：status='BANNED'；解封：status='ACTIVE'
    \item 输出：成功消息
    \item 权限：ADMIN
\end{itemize}

\textbf{③ 商品审核（M6-F04/F05）}

\begin{itemize}
    \item 输入：page, status筛选
    \item 处理：分页查询products表 → 审核通过：status='ACTIVE'；下架：status='BANNED'
    \item 输出：Page<ProductVO>
    \item 权限：ADMIN
\end{itemize}

\textbf{④ 分类管理（M6-F06）}

\begin{itemize}
    \item 输入：CRUD操作（name, icon, parent\_id）
    \item 处理：增删改查categories表（树形结构，parent\_id关联实现无限层级）
    \item 输出：全部分类列表（树形JSON）
    \item 权限：ADMIN
\end{itemize}

\textbf{⑤ 数据概览（M6-F07）}

\begin{itemize}
    \item 输入：无
    \item 处理：聚合查询 --- COUNT(users)总用户数、COUNT(products)总商品数、今日订单数、今日成交额、拍卖成交率
    \item 输出：StatsOverviewVO
    \item 权限：ADMIN
\end{itemize}

\textbf{⑥ 趋势数据（M6-F08）}

\begin{itemize}
    \item 输入：days（7或30天）
    \item 处理：查询近N天每日新增用户数/新增订单数/成交额，按日期分组
    \item 输出：List<TrendDataVO>
    \item 权限：ADMIN
\end{itemize}

\subsubsection{相关功能}

\begin{center}
\renewcommand{\arraystretch}{1.3}
\begin{tabular}{|p{3cm}|p{4cm}|p{7cm}|}
\hline
\rowcolor{tableheader}
\textbf{组件名} & \textbf{说明} & \textbf{关系} \\
\hline
M1 用户认证 & 用户管理与权限 & ADMIN角色校验；封禁后用户Token立即失效 \\
\hline
M2 商品管理 & 商品审核 & 审核通过/下架商品 \\
\hline
M3 拍卖引擎 & 拍卖监管 & 查看进行中拍卖、异常出价监控 \\
\hline
M5 订单交易 & 交易统计 & 聚合订单数据生成统计图表 \\
\hline
\end{tabular}
\end{center}

\subsubsection{用户接口}

\begin{center}
\renewcommand{\arraystretch}{1.3}
\small
\begin{tabular}{|c|c|c|p{5cm}|}
\hline
\rowcolor{tableheader}
\textbf{接口} & \textbf{方法} & \textbf{权限} & \textbf{说明} \\
\hline
\verb|/api/v1/admin/users| & GET & ADMIN & 用户列表（分页+搜索+状态筛选） \\
\hline
\verb|/api/v1/admin/users/{id}/verify| & PUT & ADMIN & 身份审核 \\
\hline
\verb|/api/v1/admin/users/{id}/ban| & PUT & ADMIN & 封禁/解封 \\
\hline
\verb|/api/v1/admin/products| & GET & ADMIN & 商品列表（含审核筛选） \\
\hline
\verb|/api/v1/admin/products/{id}/audit| & PUT & ADMIN & 商品审核/下架 \\
\hline
\verb|/api/v1/admin/categories| & GET/POST/PUT/DELETE & ADMIN & 分类管理 \\
\hline
\verb|/api/v1/admin/stats/overview| & GET & ADMIN & 数据概览 \\
\hline
\verb|/api/v1/admin/stats/trend| & GET & ADMIN & 趋势数据 \\
\hline
\end{tabular}
\end{center}

% #################### 第5章 非功能需求 ####################
\section{非功能需求}

在非功能需求中，要明确以下内容。

\begin{itemize}
    \item 在接受需求定义书中的非功能要求时，应该明确本功能有哪些需求
    \item 在基本设计阶段，对于已明确的需求，必要时对影响范围进行讨论
\end{itemize}

\subsection{权限需求}

使用图、表的方式，明确以下内容的对应关系：

\begin{itemize}
    \item 用户、用户权限、功能模块
    \item 其他特殊权限需求
\end{itemize}

\begin{center}
\renewcommand{\arraystretch}{1.3}
\begin{tabular}{|c|c|p{3cm}|p{6cm}|}
\hline
\rowcolor{tableheader}
\textbf{权限级别} & \textbf{角色标识} & \textbf{说明} & \textbf{包含权限} \\
\hline
未登录 & ANONYMOUS & 未登录访客 & 浏览商品列表、查看商品详情、搜索商品、查看用户公开主页 \\
\hline
普通用户 & USER & 已认证用户 & 以上全部 + 发布商品、出价、IM聊天、下单支付、确认收货、评价交易、管理个人信息 \\
\hline
管理员 & ADMIN & 平台运营维护者 & 以上全部 + 用户管理(审核/封禁)、商品管理(审核/下架)、分类管理、拍卖监管、数据统计 \\
\hline
\end{tabular}
\end{center}

\textbf{实现方式}：Spring Security + JWT + \verb|@PreAuthorize|注解。角色存储于users表的role字段。Service层通过SecurityContextHolder.getContext()获取当前用户进行数据归属校验。

\subsection{兼容互换性}

必须对与过去的版本有兼容性或与相关产品有兼容性的事情进行关注并要明确记述。

\begin{itemize}
    \item \textbf{浏览器兼容}：桌面端 Chrome 120+ / Edge 120+ / Firefox 120+；移动端 iOS Safari 16+ / Android Chrome 120+
    \item \textbf{API版本兼容}：使用 \verb|/api/v1/| 路径前缀，后续版本可并行运行
    \item \textbf{数据库兼容}：标准MySQL 8.0 SQL，可迁移至MariaDB 10.6+或云数据库RDS
    \item \textbf{容器兼容}：Docker Compose编排，任何支持Docker的环境均可运行
    \item \textbf{图片格式兼容}：JPEG、PNG、GIF、WebP
\end{itemize}

\subsection{安全性}

以下内容是功能性中所存在的重要的副特性。

\begin{itemize}
    \item 标准的符合性（是否符合基准，标准）
    \item 安全性（是否能够防止信息的泄露，丢失）
    \item 给出用户在安全方面的需求
\end{itemize}

\begin{center}
\renewcommand{\arraystretch}{1.3}
\begin{tabular}{|p{4cm}|p{10cm}|}
\hline
\rowcolor{tableheader}
\textbf{安全措施} & \textbf{实现方式} \\
\hline
密码加密 & BCryptPasswordEncoder（强度10），不可逆 \\
\hline
身份认证 & JWT HMAC-SHA256签名，密钥≥256位。AccessToken 30分钟，RefreshToken 7天 \\
\hline
接口授权 & @PreAuthorize注解 + 自定义权限表达式。Service层校验数据归属 \\
\hline
SQL注入防护 & MyBatis-Plus \verb|#{}|参数化查询（100\%防注入），严禁使用 \verb|${}| \\
\hline
XSS防护 & 输入：Spring HtmlUtils.htmlEscape()过滤；输出：Vue自动转义 \\
\hline
CSRF & 前后端分离（Token认证而非Cookie认证），天然免疫 \\
\hline
文件上传安全 & 类型白名单(jpg/png/gif/webp)、大小≤5MB、UUID随机文件名 \\
\hline
CORS & 仅允许前端域名白名单 + localhost:5173 \\
\hline
IM内容审核 & 本地敏感词词库（sensitive\_words表），启动时加载到内存HashSet \\
\hline
端口隔离 & Docker仅暴露80/443（Nginx），MySQL和Redis仅容器内网访问 \\
\hline
\end{tabular}
\end{center}

\subsection{健壮性}

明确本制品所要求达到的健壮性。

\begin{itemize}
    \item 成熟性（故障的频率是怎样的）
    \item 障害兼容性（在发生障害时，能够兼容到一个什么样的程度）
    \item 恢复性（发生故障时的恢复能力）
\end{itemize}

\begin{center}
\renewcommand{\arraystretch}{1.3}
\begin{tabular}{|p{4cm}|p{10cm}|}
\hline
\rowcolor{tableheader}
\textbf{措施} & \textbf{说明} \\
\hline
统一异常处理 & @RestControllerAdvice + @ExceptionHandler，所有未捕获异常返回统一格式 \\
\hline
输入校验 & 前端Element Plus表单验证（失焦+提交）；后端Spring Validation（@Valid/@NotBlank等） \\
\hline
WebSocket断连重连 & 前端SockJS+STOMP，指数退避 1s→2s→4s→8s→16s→最大30s \\
\hline
数据库连接池 & HikariCP，最大连接数20，最小空闲5，连接超时30s，泄漏检测60s \\
\hline
优雅关闭 & server.shutdown=graceful，等待进行中请求完成（最多30秒） \\
\hline
定时任务幂等 & @Scheduled拍卖判定使用Redis分布式锁（SETNX auction:check:lock NX EX 60） \\
\hline
文件上传异常隔离 & 单张图片上传失败不影响同批次其他图片 \\
\hline
\end{tabular}
\end{center}

\subsection{使用性（操作性）}

明确本制品所要求达到使用性。

\begin{itemize}
    \item 记述所要实现功能的用户界面，操作性等
    \item 需要在线帮助时，要定义其基本内容
\end{itemize}

\begin{center}
\renewcommand{\arraystretch}{1.3}
\begin{tabular}{|p{4cm}|p{10cm}|}
\hline
\rowcolor{tableheader}
\textbf{措施} & \textbf{说明} \\
\hline
响应式布局 & CSS Flexbox/Grid，适配1920px桌面至375px手机 \\
\hline
加载状态 & 骨架屏（首页商品列表）、Loading Spinner（按钮）、进度条（图片上传） \\
\hline
操作反馈 & ElMessage Toast：成功绿色 / 警告橙色 / 错误红色，3秒自动消失 \\
\hline
防重复提交 & 出价按钮出价中锁定1秒；下单按钮变"处理中..." \\
\hline
空状态 & 列表无数据显示插画+"暂无数据"文字 \\
\hline
返回顶部 & 滚动超过一屏显示浮动按钮 \\
\hline
草稿保存 & 发布商品表单内容localStorage暂存 \\
\hline
快捷键 & 聊天窗口 Enter发送、Shift+Enter换行 \\
\hline
\end{tabular}
\end{center}

\subsection{效率性（性能）}

明确所要实现功能的界限值和要求的动作。

\textbf{该项目用于性能测试。}

\begin{center}
\renewcommand{\arraystretch}{1.3}
\begin{tabular}{|c|c|p{7cm}|}
\hline
\rowcolor{tableheader}
\textbf{指标} & \textbf{目标值} & \textbf{实现策略} \\
\hline
核心API P50延迟 & < 200ms & HikariCP连接池 + MySQL索引覆盖高频查询 \\
\hline
核心API P99延迟 & < 500ms & Redis缓存热点数据（商品列表5min、用户信息10min） \\
\hline
拍卖出价推送延迟 & < 500ms & Redis Lua原子(<1ms) + STOMP直推 \\
\hline
商品搜索响应 & < 300ms & MySQL ngram全文索引 + LIKE降级 \\
\hline
并发用户数 & ≥ 200 & Tomcat max-threads=200；无状态设计支持水平扩展 \\
\hline
系统吞吐量 & ≥ 300 TPS & 拍卖出价异步持久化(@Async)不阻塞主线程 \\
\hline
\end{tabular}
\end{center}

\textbf{数据库索引清单：}

\begin{lstlisting}[language=sql]
-- 商品
CREATE INDEX idx_products_status_created ON products(status, created_at DESC);
CREATE FULLTEXT INDEX ft_products_search ON products(title, description)
    WITH PARSER ngram;
-- 拍卖
CREATE INDEX idx_bids_product_time ON bids(product_id, bid_time DESC);
-- IM
CREATE INDEX idx_messages_conv_time ON messages(conversation_id, created_at);
-- 订单
CREATE INDEX idx_orders_buyer_status ON orders(buyer_id, status);
CREATE INDEX idx_orders_seller_status ON orders(seller_id, status);
-- 通知
CREATE INDEX idx_notif_user_read ON notifications(user_id, is_read, created_at DESC);
\end{lstlisting}

\subsection{维护性}

明确功能所要求的维护性。例如，版本升级应该如何实施，如何提高分析性等问题明确化。

\begin{center}
\renewcommand{\arraystretch}{1.3}
\begin{tabular}{|p{3cm}|p{11cm}|}
\hline
\rowcolor{tableheader}
\textbf{措施} & \textbf{说明} \\
\hline
代码规范 & 后端 阿里巴巴Java开发手册；前端 Vue 3官方风格指南 + ESLint + Prettier \\
\hline
日志 & SLF4J+Logback。INFO：登录/出价/下单/支付；WARN：业务异常；ERROR：系统异常含堆栈 \\
\hline
数据库迁移 & Flyway版本化，脚本命名 \verb|V{版本号}__{描述}.sql| \\
\hline
API文档 & SpringDoc OpenAPI自动生成 \\
\hline
模块解耦 & 6个模块按功能分包，模块间通过Service接口调用 \\
\hline
代码注释 & 所有public方法Javadoc；复杂逻辑详细行内注释 \\
\hline
\end{tabular}
\end{center}

\subsection{移植性}

对向其他OS移植，通用化的可能性和必要性等问题进行讨论。

\begin{itemize}
    \item Docker Compose编排全部服务，\verb|docker-compose up -d| 即可在任意Linux服务器运行
    \item 不依赖特定云平台（阿里云/腾讯云/华为云/AWS均可）
    \item 标准MySQL 8.0 SQL，可迁至MariaDB 10.6+或云RDS
    \item 图片存储可切换：\verb|FileStorageService|接口抽象，默认本地文件系统，可换OSS/MinIO
\end{itemize}

\subsection{用户文档}

列举出将与软件一同发行的用户文档部分，例如，用户手册，在线帮助和教程。

\begin{itemize}
    \item README.md --- 环境要求、启动步骤、Docker部署、项目结构
    \item Swagger API文档 --- 全部接口描述+在线调试
    \item 项目总结PPT --- 演示用
\end{itemize}

\subsection{其他}

虽然是在功能基本要求或上述的质量特性中没有涉及到的事项，但作为产品来说，重要的事情应该在这里明确化。

无额外非功能需求。

% #################### 第6章 使用和操作方法 ####################
\section{使用和操作方法}

\subsection{环境设定}

描述执行本功能时所需环境的设定。

\textbf{【本地开发环境搭建】}

\textbf{Step 1 --- 启动基础设施：}

\begin{lstlisting}[language=bash]
cd campus-auction/deploy
docker-compose up -d mysql redis
\end{lstlisting}

\textbf{Step 2 --- 启动后端（Flyway首次启动自动建表）：}

\begin{lstlisting}[language=bash]
cd campus-auction/backend
./mvnw spring-boot:run
# 验证: http://localhost:8080/swagger-ui.html
\end{lstlisting}

\textbf{Step 3 --- 启动前端：}

\begin{lstlisting}[language=bash]
cd campus-auction/frontend
npm install
npm run dev
# 验证: http://localhost:5173
\end{lstlisting}

\textbf{【生产环境一键部署】}

\begin{lstlisting}[language=bash]
cd campus-auction/deploy
docker-compose up -d --build
# 启动顺序: MySQL → Redis → Spring Boot → Nginx
\end{lstlisting}

\subsection{使用方法}

描述执行本系统的方法。包括本系统使用环境，安装，和启动，退出。

\textbf{【核心用户旅程】}

\textbf{1. 注册与登录}：访问首页 → 右上角"注册" → 填写学号/邮箱/密码/手机号 → 提交 → "注册成功，等待审核" → 管理员后台审核通过 → 登录 → 获得JWT → 进入首页

\textbf{2. 浏览与搜索商品}：首页Feed流浏览 → 点击分类标签筛选 → 搜索框输入关键词 → 搜索结果页按价格/时间排序 → 点击商品卡片进入详情页

\textbf{3. 发布商品}：点击"发布" → 填写标题描述 → 拖拽上传图片 → 选择分类 → 选择交易模式（一口价/拍卖）→ 提交 → 审核通过后公开展示

\textbf{4. 参与拍卖}：进入拍卖商品详情页 → 查看当前最高价和出价记录 → 输入出价金额 → 点击"出价"→ 成功则页面实时更新 → 倒计时归零 → 系统判定得主

\textbf{5. IM私聊}：商品详情页"联系卖家"→ 进入聊天窗口 → 输入文字 → 消息实时送达 → 可发送商品实物图

\textbf{6. 订单与评价}：拍卖成交/购买后 → "我的订单"→ 点击"模拟支付"→ 卖家发货后→"确认收货"→ 弹出评价窗口→ 打分+文字→提交→双方信誉分更新

\textbf{7. 管理后台}：管理员登录 → "管理后台"入口 → 仪表盘概览 → "用户管理"审核新用户 → "商品管理"审核商品 → "数据统计"查看趋势图表

% #################### 第7章 注意限制事项 ####################
\section{注意限制事项}

\subsection{制约/限制事项}

本规格书中记述的某些功能在实现时，要明确设计上有制约限制的项目。

\begin{center}
\renewcommand{\arraystretch}{1.3}
\begin{tabular}{|p{4cm}|p{10cm}|}
\hline
\rowcolor{tableheader}
\textbf{限制项} & \textbf{说明} \\
\hline
支付限制 & 使用模拟支付（点击按钮即视为支付成功），不接入微信支付、支付宝等真实支付网关 \\
\hline
通知限制 & 使用站内通知（notifications表+WebSocket）代替短信和邮件 \\
\hline
存储限制 & 开发/演示阶段使用本地文件系统，生产建议切换OSS/MinIO \\
\hline
移动端限制 & 仅Web响应式适配（浏览器访问），不提供iOS/Android原生App \\
\hline
推荐限制 & 使用"热门排行"简单排序，不实现推荐算法 \\
\hline
认证限制 & 学号+管理员人工审核，不接入学校SSO/CAS统一认证 \\
\hline
时间基准 & 拍卖以服务器时间判定到期，前端倒计时仅展示用 \\
\hline
并发限制 & 设计支持200并发，超此规模需增加负载均衡和数据库读写分离 \\
\hline
IM限制 & 仅支持文字和图片消息，不支持视频/文件/语音 \\
\hline
商品分页 & 列表分页固定20条/页 \\
\hline
\end{tabular}
\end{center}

\subsection{假定事项}

除了本产品所要求的运行环境（参照2.2运行环境）以外，在功能实现上，要明确假定的项目。

\begin{enumerate}
    \item 用户通过学号/工号注册即视为校园身份有效，管理员审核过后账号激活
    \item 服务器运行在稳定网络环境（可用性≥99.5\%，带宽≥100Mbps）
    \item 并发用户量在200以内（课程演示规模）
    \item 用户浏览器支持WebSocket（Chrome/Edge/Firefox 120+均支持）
    \item 服务器磁盘≥40GB（商品图片存储）
    \item 数据量在10万条以内时MySQL单实例满足性能需求
    \item 服务器系统时间与NTP同步，误差<1秒
    \item 以单实例Docker部署（非集群），拍卖定时任务通过Redis锁防多实例并发
\end{enumerate}

\subsection{特记事项}

在开发作业时，对有关该产品应该注意事项/制约方面，对于开发人员应该了解的事项和对顾客应该说明的事项进行记述。

\begin{enumerate}
    \item \textbf{拍卖时间以服务器为准}：前端倒计时以\verb|auction_end - serverNow|为基准，首次加载从服务器获取当前时间戳校准
    \item \textbf{拍卖成交不可逆}：成交后自动生成订单，买卖双方不可取消（与一口价订单不同）
    \item \textbf{出价防抖}：同一用户3秒内对同一商品多次出价，仅接受最后一次
    \item \textbf{敏感词过滤使用本地词库}：sensitive\_words表，启动时加载到内存，不调第三方API
    \item \textbf{信誉分=100+好评×1-差评×2}，最低0分。信誉分<60限制发布商品
    \item \textbf{订单默认好评}：完成后7天未评价，系统自动生成5星好评
    \item \textbf{IM会话限制}：每用户最多10个活跃会话
    \item \textbf{用户举报}：核实3次自动封禁7天
\end{enumerate}

% #################### 第8章 附录 ####################
\section{附录}

\subsection{业务规则}

业内的业务规则，如果有的话，请列出。

\begin{center}
\renewcommand{\arraystretch}{1.3}
\begin{tabular}{|c|c|p{10cm}|}
\hline
\rowcolor{tableheader}
\textbf{编号} & \textbf{规则描述} \\
\hline
BR-01 & 出价金额必须大于当前最高价，最小加价幅度1元 \\
\hline
BR-02 & 拍卖到期，最高出价者得；无出价则流拍（status=EXPIRED） \\
\hline
BR-03 & 一口价商品被购买后立即status=SOLD，不再展示 \\
\hline
BR-04 & 订单仅在PENDING状态可由买家取消；PAID后不可取消 \\
\hline
BR-05 & 交易完成后7天内双方可评价，超期默认好评 \\
\hline
BR-06 & 信誉分=100+好评×1-差评×2，<60限制发布商品 \\
\hline
BR-07 & 同一用户3秒内同一商品多次出价，仅记录最后一次 \\
\hline
BR-08 & 每用户最多10个IM会话 \\
\hline
BR-09 & 新用户首次发布商品须管理员审核；信誉分≥80免审 \\
\hline
BR-10 & 用户被举报核实3次，自动封禁7天 \\
\hline
\end{tabular}
\end{center}

\subsection{术语表}

行业术语，请列出。

\begin{center}
\renewcommand{\arraystretch}{1.3}
\begin{tabular}{|p{4cm}|p{10cm}|}
\hline
\rowcolor{tableheader}
\textbf{术语} & \textbf{定义} \\
\hline
CAIP & Campus Auction \& IM Platform，本系统英文简称 \\
\hline
英式拍卖 & 出价从低到高的升价拍卖，在截止时间内价高者得 \\
\hline
一口价 & 卖家设定固定价格，买家直接购买，无需竞价 \\
\hline
流拍 & 拍卖结束时无人出价，商品状态变为EXPIRED \\
\hline
STOMP & Simple Text Oriented Messaging Protocol，面向文本的消息协议 \\
\hline
JWT & 无状态认证令牌，包含用户ID、角色、过期时间，HMAC-SHA256签名 \\
\hline
MyBatis-Plus & Java持久层框架，提供Lambda表达式查询、自动分页、代码生成器 \\
\hline
Flyway & 数据库版本迁移工具，通过SQL脚本管理数据库Schema变更 \\
\hline
DFD & 数据流图，描述系统中数据的流动方向和处理逻辑 \\
\hline
RBAC & 基于角色的访问控制，通过角色控制用户权限 \\
\hline
HikariCP & 高性能JDBC数据库连接池 \\
\hline
软删除 & 标记deleted\_at字段而非物理删除数据 \\
\hline
信誉分 & 用户交易信用的量化指标，100+好评×1-差评×2 \\
\hline
\end{tabular}
\end{center}

\subsection{数据流图}

画出系统的数据流图。

\begin{figure}[H]
\centering
{\small
\begin{verbatim}
                    顶层数据流图 DFD Level 0

  买家 ──→ 浏览/出价/下单/支付/IM/评价  ──→
                                           ┌──────────┐
  买家 ←── 商品列表/出价结果/订单状态/IM  ←─│          │
                                           │  CAIP    │←──→ MySQL 8.0 (D1)
  卖家 ──→ 发布商品/发货/IM/评价      ──→ │  系统    │←──→ Redis 7  (D2)
                                           │          │──→  文件系统  (D3)
  卖家 ←── 成交通知/订单通知/IM        ←──│          │
                                           └──────────┘
  管理员 ──→ 管理用户/审核商品/查看统计 ──→
  管理员 ←── 审核结果/统计图表          ←──
\end{verbatim}
}
\caption{顶层数据流图 DFD Level 0}
\label{fig:dfd0}
\end{figure}

\begin{figure}[H]
\centering
{\small
\begin{verbatim}
                    一级数据流图 DFD Level 1

  买家 ──→ P1 商品浏览与搜索 ←──→ D1 products
  买家 ──→ P2 拍卖竞价引擎   ←──→ D2 bids
  买家 ──→ P3 IM即时通讯      ←──→ D3 messages
  买家 ──→ P4 订单交易管理    ←──→ D4 orders
  买家 ──→ P5 评价信誉系统    ←──→ D5 reviews

  卖家 ──→ P1 (发布商品)
  卖家 ──→ P3 (发送消息)
  卖家 ──→ P4 (发货)
  卖家 ──→ P5 (评价)

  管理员 ──→ P6 用户管理      ←──→ D6 users
  管理员 ──→ P7 商品审核      ←──→ D1 products
  管理员 ──→ P8 数据统计      ←──(只读) D1, D4
\end{verbatim}
}
\caption{一级数据流图 DFD Level 1}
\label{fig:dfd1}
\end{figure}

\subsection{数据字典}

使用表格的方式列出系统中所有对象对应的数据字典。对象、属性、说明、取值范围，特别说明（备考）。

\begin{figure}[H]
\centering
{\small
\begin{verbatim}
                          数据库ER图

  ┌──────────┐         ┌──────────────┐         ┌──────────┐
  │  users   │1──────*│   products   │1──────*│   bids    │
  │          │ (发布)  │              │ (被出价) │          │
  │ id PK    │         │ id PK        │         │ id PK    │
  │student_id│         │ seller_id FK │         │product_id│
  │ username │         │ title        │         │ user_id  │
  │ pass_hash│         │ type         │         │bid_amount│
  │ email    │         │ start_price  │         │ bid_time │
  │ role     │         │current_price │         └──────────┘
  │ rep_score│         │ status       │
  │ status   │         │ auction_end  │
  └──────────┘         │ winner_id FK │
       │               └──────────────┘
       │                      │
       │ 1              1     │
       ├──────*         ├─────*
       │               │
       ▼               ▼
  ┌──────────┐   ┌──────────────┐   ┌──────────────┐
  │ messages │   │    orders    │   │  reviews     │
  │ id PK    │   │ id PK        │   │ id PK        │
  │ conv_id  │   │ product_id FK│   │ order_id FK  │
  │sender_id │   │ buyer_id FK  │   │reviewer_id FK│
  │ content  │   │ seller_id FK │   │target_user_id│
  │cont_type │   │ amount       │   │ rating       │
  │ is_read  │   │ status       │   │ content      │
  └──────────┘   └──────────────┘   └──────────────┘
       │
       │ *
  ┌────┴─────────┐
  │conversations │
  │ id PK        │
  │ user1_id FK  │
  │ user2_id FK  │
  │ product_id FK│
  └──────────────┘
\end{verbatim}
}
\caption{数据库ER图}
\label{fig:er}
\end{figure}

\subsubsection{各表详细字段定义}

\textbf{(1) users 用户表}

\begin{center}
\renewcommand{\arraystretch}{1.3}
\small
\begin{tabular}{|c|c|c|p{6cm}|}
\hline
\rowcolor{tableheader}
\textbf{字段} & \textbf{类型} & \textbf{约束} & \textbf{说明} \\
\hline
id & BIGINT & PK, AUTO\_INCREMENT & 主键 \\
\hline
student\_id & VARCHAR(20) & UNIQUE, NOT NULL & 学号/工号 \\
\hline
username & VARCHAR(50) & NOT NULL & 用户名（昵称） \\
\hline
password\_hash & VARCHAR(255) & NOT NULL & BCrypt加密密码 \\
\hline
email & VARCHAR(100) & NOT NULL & 邮箱 \\
\hline
phone & VARCHAR(20) & NULL & 手机号 \\
\hline
avatar & VARCHAR(500) & NULL & 头像URL \\
\hline
role & VARCHAR(20) & NOT NULL, DEFAULT 'USER' & USER / ADMIN \\
\hline
reputation\_score & INT & DEFAULT 100 & 信誉分 \\
\hline
status & VARCHAR(20) & NOT NULL, DEFAULT 'PENDING' & ACTIVE / BANNED / PENDING \\
\hline
created\_at & DATETIME & NOT NULL, DEFAULT NOW() & 创建时间 \\
\hline
updated\_at & DATETIME & NOT NULL, DEFAULT NOW() ON UPDATE NOW() & 更新时间 \\
\hline
\end{tabular}
\end{center}

索引：uk\_users\_student\_id UNIQUE(student\_id), idx\_users\_status(status)

\textbf{(2) products 商品表}

\begin{center}
\renewcommand{\arraystretch}{1.3}
\small
\begin{tabular}{|c|c|c|p{6cm}|}
\hline
\rowcolor{tableheader}
\textbf{字段} & \textbf{类型} & \textbf{约束} & \textbf{说明} \\
\hline
id & BIGINT & PK, AUTO\_INCREMENT & 主键 \\
\hline
seller\_id & BIGINT & FK→users.id, NOT NULL & 卖家ID \\
\hline
title & VARCHAR(200) & NOT NULL & 标题 \\
\hline
description & TEXT & NULL & 描述 \\
\hline
category\_id & INT & FK→categories.id, NULL & 分类ID \\
\hline
type & VARCHAR(10) & NOT NULL & SALE / AUCTION \\
\hline
start\_price & DECIMAL(10,2) & NULL & 起拍价（AUCTION） \\
\hline
fixed\_price & DECIMAL(10,2) & NULL & 一口价（SALE） \\
\hline
current\_price & DECIMAL(10,2) & NULL & 当前最高价（AUCTION） \\
\hline
status & VARCHAR(20) & NOT NULL, DEFAULT 'ACTIVE' & ACTIVE / SOLD / EXPIRED / BANNED / PENDING \\
\hline
auction\_start & DATETIME & NULL & 拍卖开始时间 \\
\hline
auction\_end & DATETIME & NULL & 拍卖结束时间 \\
\hline
winner\_id & BIGINT & FK→users.id, NULL & 得主ID \\
\hline
view\_count & INT & DEFAULT 0 & 浏览次数 \\
\hline
created\_at & DATETIME & NOT NULL & 创建时间 \\
\hline
updated\_at & DATETIME & NOT NULL & 更新时间 \\
\hline
\end{tabular}
\end{center}

索引：idx\_products\_status\_created(status, created\_at DESC), idx\_products\_category(category\_id), idx\_products\_seller(seller\_id), FULLTEXT ft\_products\_search(title, description) WITH PARSER ngram

\textbf{(3) product\_images 商品图片表}

\begin{center}
\renewcommand{\arraystretch}{1.3}
\begin{tabular}{|c|c|c|p{5cm}|}
\hline
\rowcolor{tableheader}
\textbf{字段} & \textbf{类型} & \textbf{约束} & \textbf{说明} \\
\hline
id & BIGINT & PK, AUTO\_INCREMENT & 主键 \\
\hline
product\_id & BIGINT & FK→products.id, NOT NULL & 商品ID \\
\hline
image\_url & VARCHAR(500) & NOT NULL & 图片URL \\
\hline
sort\_order & INT & DEFAULT 0 & 排序 \\
\hline
\end{tabular}
\end{center}

\textbf{(4) categories 分类表}

\begin{center}
\renewcommand{\arraystretch}{1.3}
\begin{tabular}{|c|c|c|p{5cm}|}
\hline
\rowcolor{tableheader}
\textbf{字段} & \textbf{类型} & \textbf{约束} & \textbf{说明} \\
\hline
id & INT & PK, AUTO\_INCREMENT & 主键 \\
\hline
name & VARCHAR(50) & NOT NULL & 分类名 \\
\hline
icon & VARCHAR(100) & NULL & 图标 \\
\hline
parent\_id & INT & NULL & 父分类ID（NULL=顶级，树形结构） \\
\hline
\end{tabular}
\end{center}

\textbf{(5) product\_favorites 收藏表}

\begin{center}
\renewcommand{\arraystretch}{1.3}
\begin{tabular}{|c|c|c|p{5cm}|}
\hline
\rowcolor{tableheader}
\textbf{字段} & \textbf{类型} & \textbf{约束} & \textbf{说明} \\
\hline
id & BIGINT & PK, AUTO\_INCREMENT & 主键 \\
\hline
user\_id & BIGINT & FK→users.id, NOT NULL & 用户ID \\
\hline
product\_id & BIGINT & FK→products.id, NOT NULL & 商品ID \\
\hline
created\_at & DATETIME & NOT NULL & 收藏时间 \\
\hline
\end{tabular}
\end{center}

索引：uk\_fav\_user\_product UNIQUE(user\_id, product\_id)

\textbf{(6) bids 出价记录表}

\begin{center}
\renewcommand{\arraystretch}{1.3}
\begin{tabular}{|c|c|c|p{5cm}|}
\hline
\rowcolor{tableheader}
\textbf{字段} & \textbf{类型} & \textbf{约束} & \textbf{说明} \\
\hline
id & BIGINT & PK, AUTO\_INCREMENT & 主键 \\
\hline
product\_id & BIGINT & FK→products.id, NOT NULL & 商品ID \\
\hline
user\_id & BIGINT & FK→users.id, NOT NULL & 出价人ID \\
\hline
bid\_amount & DECIMAL(10,2) & NOT NULL & 出价金额 \\
\hline
bid\_time & DATETIME & NOT NULL & 出价时间 \\
\hline
\end{tabular}
\end{center}

索引：idx\_bids\_product\_time(product\_id, bid\_time DESC), uk\_bids\_product\_user UNIQUE(product\_id, user\_id, bid\_amount)

\textbf{(7) conversations 会话表}

\begin{center}
\renewcommand{\arraystretch}{1.3}
\small
\begin{tabular}{|c|c|c|p{5cm}|}
\hline
\rowcolor{tableheader}
\textbf{字段} & \textbf{类型} & \textbf{约束} & \textbf{说明} \\
\hline
id & BIGINT & PK, AUTO\_INCREMENT & 主键 \\
\hline
user1\_id & BIGINT & FK→users.id, NOT NULL & 参与方1（较小ID） \\
\hline
user2\_id & BIGINT & FK→users.id, NOT NULL & 参与方2（较大ID） \\
\hline
product\_id & BIGINT & FK→products.id, NULL & 关联商品 \\
\hline
last\_message\_at & DATETIME & NULL & 最后消息时间 \\
\hline
created\_at & DATETIME & NOT NULL & 创建时间 \\
\hline
\end{tabular}
\end{center}

索引：uk\_conv\_users\_product UNIQUE(user1\_id, user2\_id, product\_id)

\textbf{(8) messages 消息表}

\begin{center}
\renewcommand{\arraystretch}{1.3}
\small
\begin{tabular}{|c|c|c|p{5cm}|}
\hline
\rowcolor{tableheader}
\textbf{字段} & \textbf{类型} & \textbf{约束} & \textbf{说明} \\
\hline
id & BIGINT & PK, AUTO\_INCREMENT & 主键 \\
\hline
conversation\_id & BIGINT & FK→conversations.id, NOT NULL & 会话ID \\
\hline
sender\_id & BIGINT & FK→users.id, NOT NULL & 发送者 \\
\hline
content & TEXT & NOT NULL & 内容（文字或图片URL） \\
\hline
content\_type & VARCHAR(10) & DEFAULT 'TEXT' & TEXT / IMAGE \\
\hline
is\_read & TINYINT(1) & DEFAULT 0 & 0未读/1已读 \\
\hline
created\_at & DATETIME & NOT NULL & 发送时间 \\
\hline
\end{tabular}
\end{center}

索引：idx\_msg\_conv\_time(conversation\_id, created\_at)

\textbf{(9) orders 订单表}

\begin{center}
\renewcommand{\arraystretch}{1.3}
\small
\begin{tabular}{|c|c|c|p{5cm}|}
\hline
\rowcolor{tableheader}
\textbf{字段} & \textbf{类型} & \textbf{约束} & \textbf{说明} \\
\hline
id & BIGINT & PK, AUTO\_INCREMENT & 主键 \\
\hline
product\_id & BIGINT & FK→products.id, NOT NULL & 商品ID \\
\hline
buyer\_id & BIGINT & FK→users.id, NOT NULL & 买家ID \\
\hline
seller\_id & BIGINT & FK→users.id, NOT NULL & 卖家ID \\
\hline
amount & DECIMAL(10,2) & NOT NULL & 成交金额 \\
\hline
status & VARCHAR(20) & NOT NULL, DEFAULT 'PENDING' & PENDING/PAID/SHIPPED/COMPLETED/CANCELLED \\
\hline
created\_at & DATETIME & NOT NULL & 下单时间 \\
\hline
paid\_at & DATETIME & NULL & 支付时间 \\
\hline
shipped\_at & DATETIME & NULL & 发货时间 \\
\hline
completed\_at & DATETIME & NULL & 完成时间 \\
\hline
\end{tabular}
\end{center}

索引：idx\_orders\_buyer\_status(buyer\_id, status), idx\_orders\_seller\_status(seller\_id, status), idx\_orders\_product(product\_id)

\textbf{(10) reviews 评价表}

\begin{center}
\renewcommand{\arraystretch}{1.3}
\small
\begin{tabular}{|c|c|c|p{5cm}|}
\hline
\rowcolor{tableheader}
\textbf{字段} & \textbf{类型} & \textbf{约束} & \textbf{说明} \\
\hline
id & BIGINT & PK, AUTO\_INCREMENT & 主键 \\
\hline
order\_id & BIGINT & FK→orders.id, UNIQUE, NOT NULL & 订单ID（一个订单一条评价） \\
\hline
reviewer\_id & BIGINT & FK→users.id, NOT NULL & 评价者 \\
\hline
target\_user\_id & BIGINT & FK→users.id, NOT NULL & 被评价者 \\
\hline
rating & TINYINT & NOT NULL, CHECK(1-5) & 评分（1-5星） \\
\hline
content & VARCHAR(500) & NULL & 评价内容 \\
\hline
created\_at & DATETIME & NOT NULL & 评价时间 \\
\hline
\end{tabular}
\end{center}

\textbf{(11) notifications 通知表}

\begin{center}
\renewcommand{\arraystretch}{1.3}
\small
\begin{tabular}{|c|c|c|p{5cm}|}
\hline
\rowcolor{tableheader}
\textbf{字段} & \textbf{类型} & \textbf{约束} & \textbf{说明} \\
\hline
id & BIGINT & PK, AUTO\_INCREMENT & 主键 \\
\hline
user\_id & BIGINT & FK→users.id, NOT NULL & 接收者 \\
\hline
type & VARCHAR(30) & NOT NULL & 类型 \\
\hline
title & VARCHAR(100) & NOT NULL & 标题 \\
\hline
content & VARCHAR(500) & NULL & 内容 \\
\hline
is\_read & TINYINT(1) & DEFAULT 0 & 已读标记 \\
\hline
related\_id & BIGINT & NULL & 关联对象ID \\
\hline
created\_at & DATETIME & NOT NULL & 通知时间 \\
\hline
\end{tabular}
\end{center}

索引：idx\_notif\_user\_read(user\_id, is\_read, created\_at DESC)

\vspace{2em}

\begin{center}
\rule{0.6\textwidth}{0.4pt}
\end{center}

\begin{center}
\zihao{5}
\textbf{文档版本}: v1.0 \\
\textbf{作成日期}: 2026-06-11 \\
\textbf{作成}: 第9组全员（顾思豪、刘子彦、段睿智、罗圣杰、罗涵、欧阳熠）

\vspace{1em}

本文档满足《软件工程实验总要求》中对功能和概要设计说明书的全部要求。
\end{center}

\end{document}
